\documentclass[11pt]{article}

\usepackage[a4paper,margin=1in]{geometry}
\usepackage{amsmath,amssymb,amsthm,mathtools,mathrsfs}
\usepackage{bm}
\usepackage{enumitem}
\usepackage{microtype}
\usepackage{xurl}
\usepackage{xcolor}
\usepackage[colorlinks=true,linkcolor=blue!55!black,citecolor=green!35!black,
  urlcolor=blue!60!black]{hyperref}
\usepackage[nameinlink,noabbrev]{cleveref}
\usepackage[backend=biber,style=alphabetic,sorting=nyt,maxbibnames=8,
  giveninits=true]{biblatex}

\addbibresource{references.bib}

\newtheorem{theorem}{Theorem}[section]
\newtheorem{proposition}[theorem]{Proposition}
\newtheorem{lemma}[theorem]{Lemma}
\newtheorem{corollary}[theorem]{Corollary}
\newtheorem{conjecture}[theorem]{Conjectural strengthening}
\theoremstyle{definition}
\newtheorem{definition}[theorem]{Definition}
\newtheorem{assumption}[theorem]{Assumption}
\theoremstyle{remark}
\newtheorem{remark}[theorem]{Remark}
\newtheorem*{proofremark}{Remark}

\newenvironment{proofroadmap}
  {\par\medskip\noindent\textit{Proof roadmap.}%
   \begin{list}{}{\leftmargin=1.5em\rightmargin=1.5em}%
   \item\small}
  {\end{list}\medskip}

\newcommand{\ii}{\mathrm{i}}
\newcommand{\dd}{\,\mathrm{d}}
\newcommand{\N}{\mathbb{N}}
\newcommand{\Z}{\mathbb{Z}}
\newcommand{\Complex}{\mathbb{C}}
\newcommand{\R}{\mathbb{R}}
\newcommand{\Sone}{\mathbb{S}^{1}}
\newcommand{\pair}[2]{\left\langle #1,#2\right\rangle}
\newcommand{\status}[1]{\par\smallskip\noindent\textbf{Status: #1}\par\smallskip}
\newcommand{\Gspace}{\mathcal{G}(k,\beta)}
\newcommand{\Coutpm}{\mathcal{C}^{\pm}_{\mathrm{out}}(k,\beta)}
\newcommand{\Aop}{\mathbb{A}(k,\beta)}
\newcommand{\Nrep}{\mathcal{N}_{\mathrm{rep}}(k,\beta)}
\newcommand{\Rc}{\mathcal{R}_{\mathrm{c}}(k,\beta)}
\newcommand{\Rg}{\mathcal{R}_{\mathrm{g}}(k,\beta)}
\newcommand{\Tr}{\operatorname{Tr}}
\newcommand{\ran}{\operatorname{ran}}
\newcommand{\graph}{\operatorname{graph}}
\newcommand{\dist}{\operatorname{dist}}
\newcommand{\ind}{\operatorname{ind}}
\newcommand{\diag}{\operatorname{diag}}
\newcommand{\Span}{\operatorname{span}}

\setlength{\parindent}{1.5em}
\setlength{\parskip}{0.2em}
\setlength{\emergencystretch}{2em}
\allowdisplaybreaks

\title{A M\"uller--Cauchy-Relation Framework for Guided Modes\\
in Periodic Line-Defect Waveguides\\[0.5em]
\large A theorem skeleton and proof draft}
\author{}
\date{\today}

\begin{document}
\maketitle

\begin{abstract}
We formulate the fixed-quasiperiod guided-mode problem for a two-dimensional
periodic line-defect waveguide as the self-adjoint eigenvalue problem obtained
by fixing the Bloch quasimomentum $\beta$ and solving for $k>0$ (equivalently,
for $k^2$) on an infinite strip.  For spectral parameters in a compact window
inside a common essential gap, square-integrability supplies the guided condition and
exponential transverse localization is a consequence rather than a pointwise
boundary condition.  The unbounded half-guides are represented by their outgoing
Cauchy relations, described through stable generalized Bloch modes and their
Jordan chains.  In the defect cell, a M\"uller-type layer-potential
representation is augmented by homogeneous Rayleigh fields.  These ingredients
define a continuous center--half-guide operator whose quotient kernel is intended to
be naturally isomorphic to the global guided-mode eigenspace.  This document is
a working theoretical draft: background results are separated from adaptations,
and the unproved representation and kernel-identification steps are accompanied
by concise proof roadmaps.
\end{abstract}

% Full title: Introduction
\section{Introduction}

Consider a dielectric structure that is periodic in one coordinate, extends
to infinity in the transverse coordinate, and differs from a perfect periodic
medium only in a finite-width central region.  At a fixed axial Bloch
quasimomentum, a guided mode is an eigenfunction that is square-integrable in
the transverse direction.  For the identical-background setting treated
there, the exact reduction of this infinite-domain eigenproblem to a bounded
defect region by Dirichlet-to-Neumann (DtN) operators is established in
\cite[Theorem~4.5]{Fliss2013}; the associated gap
eigenfunctions are exponentially localized by
\cite[Theorem~3.5]{Fliss2013}.  Propagation-operator and Riccati constructions
for periodic outlets were developed in \cite{JolyLiFliss2006,Coatleven2012},
and Robin-to-Robin maps provide a regular alternative at frequencies where a
Dirichlet chart fails \cite{FlissKlindworthSchmidt2015}.

The purpose of this draft is not to reprove that bounded-domain reduction.
Instead, it prepares a precise continuous framework in which the center cell is
represented by a M\"uller-type boundary integral ansatz while each half-guide is
represented by its complete Cauchy relation.  The latter is a linear relation
between Dirichlet and Neumann traces; it reduces to a DtN graph when its
Dirichlet projection is invertible.  The modal coordinates must include
generalized Floquet modes and Jordan chains.  Ordinary Bloch eigenvectors alone
need not be complete, whereas Riesz-basis and translation-operator results are
available in \cite[Theorem~2.2]{HohageSoussi2013}; related spectral
decompositions and modal DtN approximations are developed in
\cite{Zhang2021,Zhang2023}.

The center representation follows the quasiperiodic/free-space architecture
of \cite[Theorem~4 and Appendix~A]{BarnettGreengard2010}, but it is deliberately
stated on a quotient.  A physical field may be unique even when its layer
densities are not.  This distinction is essential because boundary integral
operators can exhibit representation-induced nullspaces or large nonphysical
inverse norms \cite{HiptmairMoiolaSpence2022}.  In particular, second-kind
structure alone does not imply that every small singular value represents a
guided mode.

The paper has four mathematical stages.  \Cref{sec:functional-setting}
defines the fixed-$\beta$ strip eigenproblem, records the two-ended
essential-spectrum formula as a proved theorem, and separates common
gaps from the compact gap windows used later.  \Cref{sec:cauchy-reduction}
defines weak Cauchy traces,
outgoing relations, their generalized-Bloch coordinates, and the exact
restriction--gluing reduction.  \Cref{sec:muller-coupling} fixes all
layer-potential conventions, introduces the center reconstruction and the
coupled operator, and isolates the representation nullspace.
\Cref{sec:main-results} states the quotient-safe main target and the remaining
Fredholm question.  Formal statements are followed by a \emph{proof roadmap}
only when a proof or adaptation remains to be completed.

% Full title: Functional Setting and the Guided-Mode Eigenvalue Problem
\section{Functional Setting and the Guided-Mode Eigenvalue Problem}
\label{sec:functional-setting}

\subsection{Geometry, material coefficient, and quasiperiodic spaces}

Let $d>0$ be the axial period and let
\[
  B:=\R\times(0,d)
\]
be the infinite period strip.  Choose two port locations $X^-<X^+$ and set
\[
  \mathcal C^0:=(X^-,X^+)\times(0,d),\qquad
  W^-:=(-\infty,X^-)\times(0,d),\qquad
  W^+:=(X^+,\infty)\times(0,d).
\]
The artificial ports are
\[
  \Gamma^-:=\{X^-\}\times(0,d),\qquad
  \Gamma^+:=\{X^+\}\times(0,d).
\]
Write $s^-= -1$ and $s^+=1$.  If $L^\pm>0$ is the transverse period of the
corresponding half-guide, its outwardly numbered cells and cell interfaces are
\begin{align*}
  \mathcal C_j^\pm
  &:=\bigl\{(x,y)\in B:(j-1)L^\pm<s^\pm(x-X^\pm)<jL^\pm\bigr\},
  &&j\in\N,\\
  \Gamma_j^\pm
  &:=\bigl\{(x,y)\in B:s^\pm(x-X^\pm)=jL^\pm\bigr\},
  &&j\in\N_0,
\end{align*}
so that $\Gamma_0^\pm=\Gamma^\pm$ and
$W^\pm=\bigcup_{j\geq1}\overline{\mathcal C_j^\pm}$ up to cell interfaces.

The center cell contains one bounded inclusion
$\Omega^0\Subset\mathcal C^0$.  A reference inclusion
$\Omega_1^\pm\Subset\mathcal C_1^\pm$ is repeated in the corresponding half-guide,
with
\[
  \Omega_j^\pm:=\Omega_1^\pm+s^\pm(j-1)L^\pm e_x,
  \qquad e_x=(1,0).
\]
All inclusion boundaries are assumed to be $C^{2}$ and separated from the cell
boundaries by a uniform positive distance.  The physical material interfaces
are
\[
  \Upsilon^0:=\partial\Omega^0,\qquad
  \Upsilon^\pm:=\bigcup_{j\geq1}\partial\Omega_j^\pm,\qquad
  \Upsilon:=\Upsilon^0\cup\Upsilon^-\cup\Upsilon^+.
\]
For the scalar TM model, let $\rho=n^2$ be the squared refractive index,
\begin{equation}
  \rho(x,y)=
  \begin{cases}
    (n^0)^2,&(x,y)\in\Omega^0,\\
    (n^\pm)^2,&(x,y)\in\Omega_j^\pm\text{ for some }j\geq1,\\
    1,&\text{otherwise},
  \end{cases}
  \qquad 0<\rho_-\leq \rho\leq \rho_+<\infty.
  \label{eq:material}
\end{equation}
More general positive piecewise-constant half-guide coefficients can be used without
changing the definitions below.

Fix a real Bloch quasimomentum $\beta\in[-\pi/d,\pi/d)$.  For any strip-like
domain $D$ having the two horizontal sides $y=0$ and $y=d$, define
\begin{equation}
  H^1_\beta(D):=
  \bigl\{u\in H^1(D):
  \left.u\right|_{y=d}
  =e^{\ii\beta d}\left.u\right|_{y=0}\bigr\},
  \label{eq:H1beta}
\end{equation}
where the boundary restrictions are understood as Sobolev traces.  This is a
closed subspace of $H^1(D)$ with the inherited norm.  The normal-derivative
quasiperiodicity is not imposed at the form level; it follows for functions in
the operator domain.

The trace spaces on a vertical port are defined by identifying the port with
$(0,d)$.  With
\[
  \beta_m:=\beta+\frac{2\pi m}{d},\qquad
  \psi_m(y):=d^{-1/2}e^{\ii\beta_my},\qquad m\in\Z,
\]
set, for $s\in\R$,
\begin{equation}
  H^s_\beta(\Gamma):=
  \left\{f=\sum_{m\in\Z}f_m\psi_m:
  \sum_{m\in\Z}\langle\beta_m\rangle^{2s}|f_m|^2<\infty\right\},
  \qquad \langle t\rangle=(1+t^2)^{1/2}.
  \label{eq:port-Sobolev}
\end{equation}
Thus $H^{-1/2}_\beta(\Gamma)$ is the anti-dual of
$H^{1/2}_\beta(\Gamma)$ under the extension of the $L^2(0,d)$ pairing.  We use
the product Cauchy space
\begin{equation}
  \mathcal H_\Gamma^\pm
  :=H^{1/2}_\beta(\Gamma^\pm)\times H^{-1/2}_\beta(\Gamma^\pm).
  \label{eq:Cauchy-space}
\end{equation}

\subsection{Guided modes and the \texorpdfstring{$\beta$}{beta}-formulation}

Fix $\beta\in[-\pi/d,\pi/d)$.  Following the formulation of
Fliss~\cite[(2.4)]{Fliss2013}, a \emph{guided mode} at wavenumber $k>0$ is a
nonzero field $u\in H^1(B)$ satisfying
\begin{equation}
  \left\{
  \begin{aligned}
    -\frac{1}{\rho}\Delta u&=k^2u&&\text{in }B,\\
    \left.u\right|_{y=d}
      &=e^{\ii\beta d}\left.u\right|_{y=0},\\
    \left.\partial_yu\right|_{y=d}
      &=e^{\ii\beta d}\left.\partial_yu\right|_{y=0}.
  \end{aligned}
  \right.
  \label{eq:guided-mode-Fliss}
\end{equation}
The condition $u\in H^1(B)$ includes square-integrability in the unbounded
$x$-direction.  Thus transverse confinement is part of the definition, while
pointwise or exponential decay will be a consequence in a spectral gap.

This paper uses the \emph{$\beta$-formulation}: $\beta$ is prescribed and the
wavenumber $k>0$ is sought.  More precisely, the problem is
\begin{equation}
  \boxed{\quad
  \text{find }k>0\text{ and }u\in D(A(\beta)),\ u\neq0,
  \quad A(\beta)u=k^2u.\quad}
  \tag{$E^\beta$}
  \label{eq:beta-formulation}
\end{equation}
\addtocounter{equation}{1}
Introduce the Sobolev space used in this formulation,
\begin{equation}
  H^1(\Delta,B):=\{u\in H^1(B):\Delta u\in L^2(B)\},
  \label{eq:H1-Delta}
\end{equation}
and set
\begin{equation}
  A(\beta):=-\frac{1}{\rho}\Delta,
  \qquad
  D(A(\beta)):=\left\{u\in H^1(\Delta,B):
  \begin{array}{l}
    \left.u\right|_{y=d}
      =e^{\ii\beta d}\left.u\right|_{y=0},\\[-1mm]
    \left.\partial_yu\right|_{y=d}
      =e^{\ii\beta d}\left.\partial_yu\right|_{y=0}
  \end{array}\right\}.
  \label{eq:Abeta}
\end{equation}
The closed-form argument of \cite[Section~3]{Fliss2013} shows that
$A(\beta)$ is self-adjoint and nonnegative in
\begin{equation}
  H:=L^2(B,\rho\,\dd x\dd y),
  \qquad \langle u,v\rangle_H:=\int_B\rho\,u\overline v\,\dd x\dd y.
  \label{eq:weighted-Hilbert-space}
\end{equation}

\begin{definition}[Fixed-$\beta$ guided modes]
\label{def:guided-space}
For fixed $\beta$ and $k>0$, the eigenspace of problem
\eqref{eq:beta-formulation} is
\begin{equation}
  \mathcal G(k,\beta)
  :=\ker(A(\beta)-k^2I).
  \label{eq:guided-space}
\end{equation}
Its nonzero members are precisely the guided modes in
\eqref{eq:guided-mode-Fliss}.  Their geometric multiplicity is
\begin{equation}
  m_{\mathrm g}(k,\beta):=\dim\mathcal G(k,\beta).
  \label{eq:geometric-multiplicity}
\end{equation}
\end{definition}

\begin{proposition}[Equivalent weak and transmission formulations]
\label{prop:weak-strong}
Suppose $\rho$ is given by \eqref{eq:material} and all interfaces are $C^2$.
A field $u$ belongs to $\mathcal G(k,\beta)$ if and only if
$u\in H^1_\beta(B)$ and
\begin{equation}
  \int_B\nabla u\cdot\nabla\overline v\,\dd x\dd y
  =k^2\int_B\rho\,u\overline v\,\dd x\dd y
  \qquad\forall v\in H^1_\beta(B).
  \label{eq:guided-weak-form}
\end{equation}
Every such field is locally $H^2$ away from the interfaces and satisfies
\begin{align}
  \Delta u+k^2\rho u&=0&&\text{in }B\setminus\Upsilon,\label{eq:strong-pde}\\
  [u]_{\Upsilon}&=0,&[\partial_\nu u]_{\Upsilon}&=0,\label{eq:TM-jumps}\\
  \left.u\right|_{y=d}
    &=e^{\ii\beta d}\left.u\right|_{y=0},&
  \left.\partial_yu\right|_{y=d}
    &=e^{\ii\beta d}\left.\partial_yu\right|_{y=0},
    \label{eq:strong-quasi}
\end{align}
where $\nu$ is either consistently chosen unit normal on each physical
interface and $[\cdot]$ denotes the trace from the $\nu$-side minus the trace
from the opposite side.  Conversely, a piecewise $H^2_{\mathrm{loc}}$ field
in $H^1(B)$ satisfying \eqref{eq:strong-pde}--\eqref{eq:strong-quasi} belongs
to $\mathcal G(k,\beta)$.
\end{proposition}
\status{background result.}
\begin{proofroadmap}
The equivalence of the operator equation and \eqref{eq:guided-weak-form}
follows from Green's identity in the weighted space $H$.  Global
$H^1$ regularity gives continuity of the Dirichlet trace across each smooth
interface; integration by parts on the two material sides shows that the
distributional interface source vanishes exactly when the two normal
derivatives agree.  The horizontal boundary terms cancel by quasiperiodicity,
and standard transmission regularity gives piecewise local $H^2$ regularity.
The converse is the same calculation in reverse.  Suitable background tools
are collected in \cite[Chapters~3--4]{McLean2000}.

Relevant internal notes are
\path{research/planning/muller-cauchy/notes/spaces.md} and
\path{research/planning/muller-cauchy/notes/trace-glue.md}.
The weakest acceptable form is distributional equivalence without asserting
global piecewise $H^2$ regularity at material junctions; such junctions are
excluded by the present smooth, separated geometry.
\end{proofroadmap}
\begin{proof}
We prove first the equivalence between the operator and weak formulations and
then recover the transmission conditions.

\emph{Step 1: the operator equation implies the weak equation.}
Let $u\in\mathcal G(k,\beta)$.  Then
$-\Delta u=k^2\rho u$ in $L^2(B)$ and both the Dirichlet and normal traces of
$u$ are $\beta$-quasiperiodic on the horizontal sides.  Take first
$v\in H^1_\beta(B)$ with bounded support in the $x$-direction.  Green's
identity on a bounded strip containing that support gives
\[
  \int_B\nabla u\cdot\nabla\overline v
  =k^2\int_B\rho u\overline v.
\]
Indeed, there are no vertical-end terms, and the terms at $y=0$ and $y=d$
cancel because the traces of $u$, $\partial_yu$, and $v$ carry the same
quasiperiodic phase.  Such test functions are dense in $H^1_\beta(B)$, so
continuity of both integrals proves \eqref{eq:guided-weak-form} for every
$v\in H^1_\beta(B)$.

\emph{Step 2: the weak equation implies the operator equation.}
Suppose that $u\in H^1_\beta(B)$ satisfies
\eqref{eq:guided-weak-form}.  Testing first with compactly supported smooth
functions shows in distributions that
\[
  -\Delta u=k^2\rho u.
\]
Since $\rho u\in L^2(B)$, this gives $u\in H^1(\Delta,B)$.  Its weak normal
traces on $y=0,d$ are therefore defined.  Applying Green's identity to a
boundedly supported $v\in H^1_\beta(B)$ and using the distributional equation
leaves only
\[
  \pair{\partial_yu|_{y=d}}{v|_{y=d}}
  -\pair{\partial_yu|_{y=0}}{v|_{y=0}}=0.
\]
Because $v|_{y=d}=e^{\ii\beta d}v|_{y=0}$ and the trace at $y=0$ can be
chosen arbitrarily in $H^{1/2}(\R)$, this identity is precisely
\[
  \partial_yu|_{y=d}
  =e^{\ii\beta d}\partial_yu|_{y=0}
  \quad\text{in }H^{-1/2}.
\]
The corresponding Dirichlet condition already follows from
$u\in H^1_\beta(B)$.  Hence $u\in D(A(\beta))$ and
$A(\beta)u=k^2u$, so $u\in\mathcal G(k,\beta)$.

\emph{Step 3: transmission conditions and local regularity.}
For a weak solution, the distributional equation found above holds with an
$L^2_{\mathrm{loc}}$ right-hand side.  Interior elliptic regularity gives
$u\in H^2_{\mathrm{loc}}$ in every material subdomain away from its boundary.
Since $u$ is one global $H^1$ function, its two Dirichlet traces on every
component of $\Upsilon$ agree.  Integrate the distributional equation by
parts on the two sides of one interface, using a test function supported in a
small interface neighbourhood.  The volume terms cancel by
\eqref{eq:strong-pde}, leaving
\[
  \pair{[\partial_\nu u]_{\Upsilon}}{\varphi}_{\Upsilon}=0
  \qquad\forall\varphi\in H^{1/2}(\Upsilon)
\]
locally on that component.  Thus both relations in
\eqref{eq:TM-jumps} hold, while \eqref{eq:strong-quasi} was established in
Steps 1--2.

\emph{Step 4: converse from the transmission formulation.}
Let now $u\in H^1(B)$ be piecewise $H^2_{\mathrm{loc}}$ and satisfy
\eqref{eq:strong-pde}--\eqref{eq:strong-quasi}.  For a
$\beta$-quasiperiodic test function with bounded $x$-support, integrate by
parts in the finitely many material pieces meeting its support.  The
Dirichlet traces define a global $H^1$ field, the interface terms cancel by
$[\partial_\nu u]_{\Upsilon}=0$, and the two horizontal terms cancel by
\eqref{eq:strong-quasi}.  This gives \eqref{eq:guided-weak-form}; density
extends it to all of $H^1_\beta(B)$.  The same calculation shows that the
piecewise Laplacians have no interface distribution, hence
$\Delta u=-k^2\rho u\in L^2(B)$.  Therefore $u\in D(A(\beta))$ and
$A(\beta)u=k^2u$, which completes the converse and the proof.
\end{proof}

\begin{remark}[TE polarization]
The draft develops only the TM model.  A scalar TE reduction has a
divergence-form principal part, for example
$-\nabla\cdot(\rho^{-1}\nabla u)=k^2u$, and continuity of the weighted normal
flux $\rho^{-1}\partial_\nu u$ replaces \eqref{eq:TM-jumps}.  This changes the form,
the conormal trace, and the integral-equation weights, so it is not treated as
a notational substitution here.
\end{remark}

\subsection{Essential spectrum, common gaps, and scan intervals}

For each $\star\in\{-,+\}$, extend the coefficient in $W^\star$
periodically in the $x$-direction to the whole strip $B$, and denote the
extension by $\rho^\star$.  The corresponding full periodic operator is
\begin{equation}
  A^\star(\beta):=-\frac{1}{\rho^\star}\Delta,\qquad
  D(A^\star(\beta)):=D(A(\beta)),
  \label{eq:background-operators}
\end{equation}
acting in $H^\star:=L^2(B,\rho^\star\,\dd x\dd y)$.  Thus
$A^\star(\beta)$ is the full periodic extension of the $\star$ half-guide,
not a half-guide operator with an artificial boundary condition at its port.

On the first cell $\mathcal C_1^\star$, introduce the outward cell coordinate
(with $X^\star=X^-$ for $\star=-$ and $X^\star=X^+$ for $\star=+$)
\[
  t^\star:=s^\star(x-X^\star)\in(0,L^\star),
  \qquad \partial_{t^\star}=s^\star\partial_x,
\]
so that $t^\star=0$ on $\Gamma_0^\star$ and $t^\star=L^\star$ on
$\Gamma_1^\star$.  For
\[
  \theta\in\mathbb B^\star
  :=\left[-\frac{\pi}{L^\star},\frac{\pi}{L^\star}\right),
\]
define the Floquet fiber in
$L^2(\mathcal C_1^\star,\rho^\star\,\dd x\dd y)$ by
\[
\begin{aligned}
  A^\star(\beta,\theta)u
  &:=-\frac{1}{\rho^\star}\Delta u,\\
  D(A^\star(\beta,\theta))
  :=\Bigl\{u\in H^1(\Delta,\mathcal C_1^\star):{}&
  \left.u\right|_{\Gamma_1^\star}
    =e^{\ii\theta L^\star}\left.u\right|_{\Gamma_0^\star},\\
  &\left.\partial_{t^\star}u\right|_{\Gamma_1^\star}
    =e^{\ii\theta L^\star}
      \left.\partial_{t^\star}u\right|_{\Gamma_0^\star},\\
  &\left.u\right|_{y=d}
    =e^{\ii\beta d}\left.u\right|_{y=0},\\
  &\left.\partial_yu\right|_{y=d}
    =e^{\ii\beta d}\left.\partial_yu\right|_{y=0}
  \Bigr\}.
\end{aligned}
\]
Here
$H^1(\Delta,\mathcal C_1^\star)
:=\{u\in H^1(\mathcal C_1^\star):\Delta u\in
L^2(\mathcal C_1^\star)\}$.  The value equalities are understood in
$H^{1/2}$ and the derivative equalities in the corresponding weak
$H^{-1/2}$ trace spaces.  Each fiber is self-adjoint and has compact
resolvent.  Denote its eigenvalues, repeated with multiplicity, by
$\lambda_n^\star(\beta,\theta)$, $n\geq1$.

Fliss~\cite[Proposition~3.1, especially equations~(3.3)--(3.4)]{Fliss2013}
gives the Floquet--Bloch band representation and continuity of the band
functions for this scalar periodic setting.  Applied separately to the two
periodic extensions, it gives
\begin{equation}
  \Sigma_\star(\beta)
  :=\sigma(A^\star(\beta))
  =\sigma_{\mathrm{ess}}(A^\star(\beta))
  =\bigcup_{n\geq1}
    \bigl\{\lambda_n^\star(\beta,\theta):
    \theta\in\mathbb B^\star\bigr\},
  \qquad \star\in\{-,+\}.
  \label{eq:background-bands}
\end{equation}
Fliss uses $\omega^2$ as the spectral parameter; in the present notation the
same parameter is $\lambda=k^2$.  The two spectra in
\eqref{eq:background-bands} must be computed independently when the periods,
cells, or material coefficients differ.

Write the nonempty bounded components of the individual positive-axis gaps as
\[
  G_m^-(\beta)
  :=\bigl(a_m^-(\beta),b_m^-(\beta)\bigr),
  \qquad
  G_n^+(\beta)
  :=\bigl(a_n^+(\beta),b_n^+(\beta)\bigr).
\]
Their common gaps are the nonempty intersections
\begin{equation}
\begin{aligned}
  G_{mn}(\beta)
  &:=G_m^-(\beta)\cap G_n^+(\beta)\\
  &=\left(
    \max\{a_m^-(\beta),a_n^+(\beta)\},
    \min\{b_m^-(\beta),b_n^+(\beta)\}
    \right),
  \label{eq:common-gap-intersection}
\end{aligned}
\end{equation}
which exist precisely when
\[
  \max\{a_m^-(\beta),a_n^+(\beta)\}
  <\min\{b_m^-(\beta),b_n^+(\beta)\}.
\]
Enumerate all such intervals as
$G_j(\beta)=(a_j(\beta),b_j(\beta))$,
$1\leq j\leq N(\beta)$, and assume below that $N(\beta)\geq1$.

For comparison, De~Nittis et al.~
\cite[Proposition~3.10 and Lemma~3.12]{DeNittisEtAl2021} prove
$\sigma_{\mathrm{ess}}(M)=\sigma(M_\ell)\cup\sigma(M_r)$
for a one-dimensional first-order Maxwell operator with two different
asymptotically periodic matrix weights.  Their Weyl essential spectrum is
characterized there by Zhislin sequences.  This is a structural precedent for
a two-ended periodic junction, but it does not itself prove the corresponding
statement for the present two-dimensional scalar second-order strip operator.
The proof in \Cref{app:essential-spectrum} adapts their localization mechanism
to the present strip and then uses the separate Floquet--Bloch descriptions
\eqref{eq:background-bands}.

\begin{theorem}[Essential spectrum for two different periodic half-guides]
\label{thm:two-ended-essential-spectrum}
\label{prop:asymmetric-essential-spectrum}
Under the geometry and coefficient assumptions of \eqref{eq:material}, the
Weyl essential spectrum of the junction operator is
\begin{equation}
  \sigma_{\mathrm{ess}}(A(\beta))
  =\Sigma_{\mathrm{ess}}(\beta)
  :=\Sigma_-(\beta)\cup\Sigma_+(\beta).
  \label{eq:projected-spectrum}
\end{equation}
In particular, no equality between the two periods or the two periodic
coefficients is required.
\end{theorem}

The detailed proof, including the local-compactness, Zhislin-sequence, and
one-sided periodic Weyl-sequence lemmas, is given in
\Cref{app:essential-spectrum}.

\begin{corollary}[Admissible scan intervals for fixed $\beta$]
\label{cor:admissible-scan-intervals}
\label{prop:admissible-scan-intervals}
If $k^2$ is an isolated $L^2$ eigenvalue of $A(\beta)$ lying outside this
essential spectrum, then
\[
  k^2\notin\Sigma_-(\beta)\cup\Sigma_+(\beta).
\]
Consequently, the candidate real-frequency scan set for ordinary discrete
gap modes is
\begin{equation}
  \mathcal K(\beta)
  :=\left\{k>0:
  k^2\notin\Sigma_-(\beta)\cup\Sigma_+(\beta)\right\}
  =\bigcup_{j=1}^{N(\beta)}
    \bigl(\sqrt{a_j(\beta)},\sqrt{b_j(\beta)}\bigr).
  \label{eq:admissible-k-scan}
\end{equation}
Equivalently, $k^2$ must lie in a common gap of the two full periodic
half-guide extensions.  This is a necessary scan condition and does not
assert that any particular common gap contains an eigenvalue.
\end{corollary}
\begin{proof}
By \Cref{thm:two-ended-essential-spectrum}, every isolated eigenvalue outside
the essential spectrum is excluded from both periodic spectra.  On the
positive spectral axis,
\[
  \bigl((0,\infty)\setminus\Sigma_-(\beta)\bigr)
  \cap
  \bigl((0,\infty)\setminus\Sigma_+(\beta)\bigr)
  =(0,\infty)\setminus
  \bigl(\Sigma_-(\beta)\cup\Sigma_+(\beta)\bigr),
\]
so the admissible spectral-parameter intervals are intersections of left and
right gaps, not their union.  For the pair $(m,n)$,
\eqref{eq:common-gap-intersection} and $\lambda=k^2$ give
\[
  k\in\left(
  \sqrt{\max\{a_m^-(\beta),a_n^+(\beta)\}},
  \sqrt{\min\{b_m^-(\beta),b_n^+(\beta)\}}
  \right),
\]
provided the strict endpoint inequality above holds.
\end{proof}

\begin{remark}[Numerical consequence]
For each fixed $\beta$, the practical scan proceeds as follows:
\begin{enumerate}[leftmargin=*]
  \item solve the left-cell Bloch problems while scanning
  $\theta\in\mathbb B^-$ and assemble $\Sigma_-(\beta)$;
  \item independently solve the right-cell Bloch problems while scanning
  $\theta\in\mathbb B^+$ and assemble $\Sigma_+(\beta)$;
  \item construct the two individual families of gaps;
  \item intersect the left and right gaps by
  \eqref{eq:common-gap-intersection};
  \item take square roots of the resulting $k^2$ intervals; and
  \item run the transmission-eigenvalue or smallest-singular-value scan only
  on those candidate $k$ intervals.
\end{enumerate}
The two periods $L^-$ and $L^+$, unit cells, and material coefficients may
all differ, so the two Bloch calculations cannot in general be reused.
The center defect can generate ordinary discrete $L^2$ localized modes only
at candidate points in these common gaps, but no existence in any candidate
interval is guaranteed.
\end{remark}

\begin{remark}[Embedded and threshold phenomena]
The scan set \eqref{eq:admissible-k-scan} targets ordinary discrete $L^2$
eigenvalues outside the essential spectrum, namely gap modes.  It does not
cover embedded eigenvalues, bound states in the continuum, threshold
eigenvalues, or band-edge resonances.  Such objects, if present, may satisfy
$k^2\in\Sigma_-(\beta)\cup\Sigma_+(\beta)$ and require a separate analysis.
Thus all isolated eigenvalues outside the essential spectrum must lie in the
common gaps; the statement is not a classification of every possible
eigenvalue.
\end{remark}

The former detailed coupled/decoupled resolvent argument, including the closed
forms, common resolvent point, Weyl theorem, and compactness route, has been
moved to
\path{archive/essential_spectrum_resolvent_argument.tex}.  It is retained
there as a technical proof attempt and is not included in this document.

\begin{proposition}[Discrete spectrum in a common gap]
\label{prop:gap-discrete-spectrum}
For every $j$ the set $\sigma(A(\beta))\cap G_j(\beta)$ consists only of isolated
eigenvalues of finite multiplicity.  Such eigenvalues can accumulate only at
$a_j(\beta)$ or $b_j(\beta)$.
\end{proposition}
\begin{proof}
By \Cref{thm:two-ended-essential-spectrum}, the interval $G_j(\beta)$ is
disjoint from the essential spectrum of the self-adjoint operator
$A(\beta)$.  Standard self-adjoint spectral theory,
used in the identical-background setting in
Fliss~\cite[Proposition~3.3]{Fliss2013}, therefore gives the claim.
\end{proof}

For the uniform estimates used in later sections, choose one common gap
$G_{j_*}(\beta)$ and a relatively compact open window
\begin{equation}
  I_\beta=(\lambda_-,\lambda_+),\qquad
  \overline{I_\beta}\subset G_{j_*}(\beta),\qquad
  \overline{I_\beta}\cap\Sigma_{\mathrm{ess}}(\beta)=\varnothing.
  \label{eq:strict-gap}
\end{equation}
The eigenproblem \eqref{eq:beta-formulation} itself is defined for every
$k>0$; the restriction $k^2\in I_\beta$ is imposed only in the later
Cauchy-relation analysis.  It guarantees simultaneously that any spectral
point is a discrete guided eigenvalue, that neither half-guide has a
propagating Bloch channel, and that the stable/unstable splitting is uniformly
separated.  To study an entire gap, one may cover it by such windows and treat
the band edges separately.

\begin{proposition}[Localization with two different half-guides]
\label{prop:asymmetric-localization}
If $k^2\in G_j(\beta)$ is an eigenvalue of $A(\beta)$, every associated
eigenfunction decays exponentially in both transverse directions.
\end{proposition}
\begin{proof}
The weighted-resolvent argument of
\cite[Proposition~3.3 and Theorem~3.5]{Fliss2013} applies separately at the two
ends.  More explicitly, the positive distances from $k^2$ to
$\sigma(A^-(\beta))$ and $\sigma(A^+(\beta))$ allow small exponential
conjugations $e^{-\alpha_-x}$ on $W^-$ and $e^{\alpha_+x}$ on $W^+$ while
keeping the two background resolvents bounded.  After a cutoff near each
port, the eigenvalue equation has a compactly supported right-hand side, so
these conjugated resolvent estimates give exponential weighted $H^1$ bounds
on the corresponding half-guides.  Joining the two estimates across the
bounded center cell proves two-sided exponential decay.  Pointwise decay then
follows from local elliptic regularity.
\end{proof}

Thus the informal condition $u(x,y)\to0$ as $x\to\pm\infty$ is a consequence
for gap eigenfunctions, not the functional definition of a guided mode.

% Full title: Outgoing Cauchy Relations and Bounded-Domain Reduction
\section{Outgoing Cauchy Relations and Bounded-Domain Reduction}
\label{sec:cauchy-reduction}

\subsection{Weak traces and outgoing half-guide spaces}

Let $D$ be one of $\mathcal C^0,W^-,W^+$.  For a field
$u\in H^1_\beta(D)$ satisfying $\Delta u\in L^2(D)$ piecewise, the Dirichlet
trace on a port belongs to $H^{1/2}_\beta(\Gamma)$.  Its weak outward Neumann
trace $\gamma_{N,D}u\in H^{-1/2}_\beta(\Gamma)$ is defined by
\begin{equation}
  \pair{\gamma_{N,D}u}{\varphi}_{\Gamma}
  :=\int_D\left(\nabla u\cdot\nabla\overline V-k^2\rho\,u\overline V\right)
  \dd x\dd y,
  \label{eq:weak-Neumann}
\end{equation}
where $V\in H^1_\beta(D)$ has port trace $\varphi$, vanishes on any other
finite port, and is otherwise an arbitrary bounded lifting.  The weak PDE
makes \eqref{eq:weak-Neumann} independent of the lifting.  Horizontal terms
cancel by quasiperiodicity.

The unit normals pointing \emph{out of the center cell} are
\[
  \nu^{0,-}=-e_x\quad\text{on }\Gamma^- ,\qquad
  \nu^{0,+}=e_x\quad\text{on }\Gamma^+.
\]
They are opposite to the outward normals of the adjacent half-guides.  We use
this center-oriented convention everywhere.  For $u^0$ in $\mathcal C^0$ and
$w^\pm$ in $W^\pm$, define
\begin{align}
  \Tr^{0,\pm} u^0
  &:=\bigl(\gamma_Du^0,\gamma_{N,\mathcal C^0}u^0\bigr)
  \in\mathcal H_\Gamma^\pm,\label{eq:center-trace}\\
  \Tr^{0,\pm} w^\pm
  &:=\bigl(\gamma_Dw^\pm,-\gamma_{N,W^\pm}w^\pm\bigr)
  \in\mathcal H_\Gamma^\pm.
  \label{eq:halfguide-trace}
\end{align}
Consequently, equality of the two pairs means continuity of the field and of
the derivative in the same physical direction $\nu^{0,\pm}$.

Let $H^1_{\beta,0}(W^\pm;\Gamma^\pm)$ be the subspace with zero Dirichlet
trace at the finite port.  The square-integrable half-guide solution space is
\begin{equation}
  \mathcal U^\pm_{\mathrm{out}}(k,\beta):=
  \left\{w\in H^1_\beta(W^\pm):
  \int_{W^\pm}\nabla w\cdot\nabla\overline v
  =k^2\int_{W^\pm}\rho\,w\overline v
  \ \forall v\in H^1_{\beta,0}(W^\pm;\Gamma^\pm)\right\}.
  \label{eq:halfguide-outgoing}
\end{equation}
When $k^2\in I_\beta$, this $H^1$ condition selects the stable, decaying
solutions.  No pointwise radiation condition is added.

\begin{definition}[Outgoing Cauchy relations]
\label{def:outgoing-relations}
The outgoing Cauchy relation of the $\pm$ half-guide is
\begin{equation}
  \mathcal C^\pm_{\mathrm{out}}(k,\beta)
  :=\left\{\Tr^{0,\pm} w:w\in
  \mathcal U^\pm_{\mathrm{out}}(k,\beta)\right\}
  \subset\mathcal H_\Gamma^\pm.
  \label{eq:outgoing-relation}
\end{equation}
It is a linear subspace of complete Dirichlet--Neumann pairs; it is not assumed
to be the graph of an operator.  The Dirichlet and Neumann projections are
\begin{equation}
  \pi_D^\pm(f,g):=f,\qquad \pi_N^\pm(f,g):=g,
  \qquad (f,g)\in\mathcal H_\Gamma^\pm.
  \label{eq:trace-projections}
\end{equation}
\end{definition}

\begin{lemma}[Weak gluing across an artificial port]
\label{lem:weak-gluing}
Let $u_1$ and $u_2$ be $H^1_\beta$ weak Helmholtz solutions in two adjacent
Lipschitz subdomains.  If their Dirichlet traces agree and their outward weak
Neumann traces sum to zero on the common port, then their piecewise union is an
$H^1_\beta$ weak solution on the union.  Conversely, the restrictions of a
global weak solution satisfy these two matching conditions.
\end{lemma}
\status{background result.}
\begin{proofroadmap}
The Sobolev gluing theorem gives a global $H^1$ function from equality of the
Dirichlet traces.  Apply Green's identity in the two subdomains.  The only new
distributional term is the sum of the two outward conormal traces; it vanishes
by hypothesis.  This is precisely why the sign in \eqref{eq:halfguide-trace} is
required.  The argument uses full Cauchy data and never infers uniqueness from
zero Dirichlet trace alone.

See \path{research/planning/muller-cauchy/notes/trace-glue.md} and
the obligation named ``restriction--gluing isomorphism'' in
\path{research/planning/muller-cauchy/p-proofs.md}.  The weakest
acceptable conclusion is distributional gluing; stronger local regularity may
then be recovered away from corners and material interfaces.
\end{proofroadmap}
\begin{proof}
Let $D_1$ and $D_2$ be the two subdomains and let
$\Gamma:=\partial D_1\cap\partial D_2$ denote their common port.  We prove
the two directions separately.

\emph{Step 1: Sobolev gluing.}
Set $u=u_j$ almost everywhere in $D_j$.  Equality of the Dirichlet traces
implies $u\in H^1_\beta(D_1\cup D_2)$.  For completeness, this follows
locally by flattening $\Gamma$ and computing the distributional derivatives:
the derivative of the piecewise function is the piecewise derivative plus an
interface distribution whose density is the jump
$\gamma_Du_1-\gamma_Du_2$.  The latter is zero in
$H^{1/2}_\beta(\Gamma)$, so the distributional derivatives belong to $L^2$.
The quasiperiodic trace condition on the horizontal sides is inherited from
the two pieces.

\emph{Step 2: cancellation of the conormal traces.}
Take $v\in H^1_\beta(D_1\cup D_2)$, initially with bounded support if the
union is unbounded.  The weak Green identities in the two subdomains give
\begin{align*}
  &\sum_{j=1}^2\int_{D_j}
  \left(\nabla u_j\cdot\nabla\overline v
  -k^2\rho\,u_j\overline v\right)\dd x\dd y\\
  &\qquad=
  \pair{\gamma_{N,D_1}u_1+\gamma_{N,D_2}u_2}
  {\gamma_Dv}_{\Gamma}.
\end{align*}
All other boundary terms either vanish by the imposed test space or cancel
by quasiperiodicity.  The right-hand side is zero by hypothesis.  Density of
the bounded-support tests in the relevant global test space then shows that
$u$ satisfies the weak Helmholtz equation on the union.

\emph{Step 3: the converse.}
Let $u\in H^1_\beta(D_1\cup D_2)$ be a global weak solution and set
$u_j=u|_{D_j}$.  A global $H^1$ function has a single Dirichlet trace on the
interior Lipschitz interface, hence
$\gamma_Du_1=\gamma_Du_2$.  Applying the two weak Green identities and
subtracting the global weak equation yields
\[
  \pair{\gamma_{N,D_1}u_1+\gamma_{N,D_2}u_2}{\varphi}_{\Gamma}=0
  \qquad\forall\varphi\in H^{1/2}_\beta(\Gamma),
\]
where an arbitrary $\varphi$ is realized by a bounded lifting supported near
$\Gamma$.  Thus the two outward weak Neumann traces sum to zero in
$H^{-1/2}_\beta(\Gamma)$, as claimed.
\end{proof}

Define the bounded center solution space
\begin{equation}
\begin{split}
  \mathcal B^0(k,\beta):=\bigl\{u^0\in H^1_\beta(\mathcal C^0):\ &
  u^0\text{ satisfies the center weak Helmholtz equation},\\[-0.2em]
  &\Tr^{0,-}u^0\in\mathcal C^-_{\mathrm{out}}(k,\beta),\quad
  \Tr^{0,+}u^0\in\mathcal C^+_{\mathrm{out}}(k,\beta)\bigr\},
\end{split}
\label{eq:bounded-center-space}
\end{equation}
where the center weak equation is tested against functions in
$H^1_\beta(\mathcal C^0)$ with zero Dirichlet trace on both ports.  The
coefficient $\rho$ encodes the physical transmission conditions at $\Upsilon^0$.

\begin{theorem}[Bounded trace-relation reduction]
\label{thm:bounded-reduction}
Assume $k^2\in I_\beta$ and that the half-guide relations in
\eqref{eq:outgoing-relation} are closed.  Restriction to the center cell defines
a linear isomorphism
\begin{equation}
  \operatorname{Res}^0:\mathcal G(k,\beta)
  \longrightarrow\mathcal B^0(k,\beta),
  \qquad u\longmapsto u|_{\mathcal C^0}.
  \label{eq:restriction-isomorphism}
\end{equation}
The inverse is obtained by choosing half-guide fields with the prescribed
Cauchy pairs and gluing them to the center field.  In particular,
\begin{equation}
  \dim\mathcal G(k,\beta)=\dim\mathcal B^0(k,\beta).
  \label{eq:bounded-multiplicity}
\end{equation}
\end{theorem}
\status{adaptation of a known result.}
\begin{proofroadmap}
For the forward map, restrict a global guided mode to the three subdomains;
the half-guide restrictions lie in \eqref{eq:halfguide-outgoing} and yield the two
relation conditions.  If the center restriction vanishes, its complete Cauchy
data vanish, and weak unique continuation gives a zero global field.

Conversely, choose half-guide preimages of the two relation pairs and apply
\Cref{lem:weak-gluing}.  Independence of the choices follows from uniqueness
for a half-guide field with zero complete Cauchy data.  Exponential decay is
then inherited from the half-guide spaces in the common gap.  For identical
half-guides this reduces to the DtN reduction of
\cite[Theorem~4.5]{Fliss2013}; the present asymmetric statement instead uses
the two Cauchy relations and the gluing argument independently at each end.

The internal audit is
\path{research/planning/muller-cauchy/r-draft-thm.md};
the detailed proof obligation is in
\path{research/planning/muller-cauchy/p-proofs.md}.  If closedness
or half-guide uniqueness fails, the weakest acceptable replacement is an
isomorphism after quotienting the half-guide zero-Cauchy solution space.
\end{proofroadmap}
\begin{proof}
We first record the uniqueness fact used below.  Under the geometry of
\Cref{sec:functional-setting}, a global piecewise Helmholtz solution which
vanishes on a nonempty open subset of the matrix phase vanishes everywhere.
Indeed, the matrix region obtained by removing the mutually separated compact
inclusions from the quasiperiodic strip is connected.  In that region the
coefficient is constant, so interior elliptic regularity and the Helmholtz
equation make the solution real analytic; vanishing on one open subset
therefore implies vanishing throughout the matrix region.  On every inclusion
boundary both the Dirichlet trace and the normal derivative consequently
vanish.  Extending the interior field by zero across a boundary patch and
using \Cref{lem:weak-gluing} gives a constant-coefficient Helmholtz solution
which vanishes on one side of that patch.  Analytic continuation makes it zero
inside the inclusion as well.  Repeating this argument proves the asserted
uniqueness on the whole strip.  The same argument applies to a half-guide
solution with zero complete Cauchy data after it is extended by zero across
its finite port.

\emph{Step 1: the restriction map is well defined and injective.}
Let $u\in\Gspace$.  Testing the global weak equation with zero extensions of
functions in $H^1_{\beta,0}(W^\pm;\Gamma^\pm)$ shows that
$u|_{W^\pm}\in\mathcal U^\pm_{\mathrm{out}}(k,\beta)$.  The center
restriction satisfies the center weak equation.  The converse part of
\Cref{lem:weak-gluing}, together with the sign convention in
\eqref{eq:halfguide-trace}, yields
\[
  \Tr^{0,\pm}(u|_{\mathcal C^0})
  =\Tr^{0,\pm}(u|_{W^\pm})
  \in\mathcal C^\pm_{\mathrm{out}}(k,\beta).
\]
Hence $\operatorname{Res}^0u\in\mathcal B^0(k,\beta)$.  If
$\operatorname{Res}^0u=0$, then $u$ vanishes on the open center matrix
region, and the uniqueness argument above gives $u=0$.  Thus
$\operatorname{Res}^0$ is injective.

\emph{Step 2: surjectivity by weak gluing.}
Take $u^0\in\mathcal B^0(k,\beta)$.  By the definition of the two outgoing
relations there are fields $w^\pm\in\mathcal U^\pm_{\mathrm{out}}(k,\beta)$
such that
\[
  \Tr^{0,\pm}w^\pm=\Tr^{0,\pm}u^0.
\]
The first components give equality of the Dirichlet traces.  By
\eqref{eq:center-trace}--\eqref{eq:halfguide-trace}, equality of the second
components is exactly
\[
  \gamma_{N,\mathcal C^0}u^0+\gamma_{N,W^\pm}w^\pm=0
  \quad\text{on }\Gamma^\pm.
\]
Applying \Cref{lem:weak-gluing} at both ports therefore produces a field
\[
  u=\begin{cases}
    w^-&\text{in }W^-,\\
    u^0&\text{in }\mathcal C^0,\\
    w^+&\text{in }W^+
  \end{cases}
  \in H^1_\beta(B)
\]
which satisfies the global weak eigenvalue equation.  The closed-form
characterization of $A(\beta)$ recorded in \Cref{prop:weak-strong} gives
$u\in\Gspace$, and by construction $\operatorname{Res}^0u=u^0$.

\emph{Step 3: independence of the half-guide preimages.}
If $\widetilde w^\pm$ is another choice with the same complete Cauchy pair,
then $z^\pm=w^\pm-\widetilde w^\pm$ has zero Dirichlet and center-oriented
Neumann traces at $\Gamma^\pm$.  Extend $z^\pm$ by zero through the center.
The gluing lemma makes this extension a global weak solution which vanishes
on a nonempty open set; the uniqueness argument at the start of the proof
implies $z^\pm=0$.  Hence the glued field is independent of all choices and
defines a linear inverse of $\operatorname{Res}^0$.

We have proved the isomorphism \eqref{eq:restriction-isomorphism}; equality of
dimensions in \eqref{eq:bounded-multiplicity} follows immediately.  The
assumed closedness of the two relations ensures that the two trace constraints
defining $\mathcal B^0(k,\beta)$ are closed; the algebraic restriction--gluing
argument itself uses only their defining trace-range property.
\end{proof}

\subsection{DtN charts and generalized Bloch traces}

Whenever $\pi_D^\pm|_{\Coutpm}$ is bijective onto
$H^{1/2}_\beta(\Gamma^\pm)$, define the DtN operator
\begin{equation}
  \Lambda^\pm(k,\beta):=
  \pi_N^\pm\bigl(\pi_D^\pm|_{\Coutpm}\bigr)^{-1}.
  \label{eq:DtN-definition}
\end{equation}
A failure of this inverse is a failure of the Dirichlet chart, not necessarily
of the Cauchy relation.  A Robin projection gives an alternative RtR chart as
in \cite{FlissKlindworthSchmidt2015}.

We next define the translation objects used below.  Extend the coefficient of
each half-guide periodically to a perfect two-sided strip.  Let
$\mathscr T^\pm(k,\beta)$ be the Hilbert space of center-oriented Cauchy traces
on $\Gamma_0^\pm$ of $H^1_{\mathrm{loc},\beta}$ solutions in that perfect
strip, equipped with a fixed local solution norm.  Unique continuation makes
the one-cell trace translation well defined on this solution-trace space:
\begin{equation}
  T^\pm(k,\beta)e:=\Tr^{0,\pm} v\big|_{\Gamma_1^\pm},
  \label{eq:translation-operator}
\end{equation}
where $v$ is the half-guide solution having trace $e$ on $\Gamma_0^\pm$, and the
trace at $\Gamma_1^\pm$ is transported back to the reference port by periodic
identification.  Boundedness of $T^\pm$ and equivalence of the local trace norm
with the solution norm are part of the translation-framework hypotheses.

\begin{lemma}[A common essential gap excludes propagating multipliers]
\label{lem:gap-unit-circle}
Assume the perfect-guide Floquet transform in the transverse periodic direction
is valid and that nonzero solution traces correspond to nonzero Bloch
solutions.  If $k^2\in I_\beta$, then
\begin{equation}
  \sigma(T^\pm(k,\beta))\cap\Sone=\varnothing.
  \label{eq:no-unit-multiplier}
\end{equation}
\end{lemma}
\status{background Floquet consequence.}
\begin{proof}
If a unit-modulus multiplier existed, write it as
$\lambda=e^{\ii\theta L^\pm}$ with real transverse quasimomentum $\theta$.
The corresponding nonzero solution would satisfy the one-cell transverse
Bloch condition and would place $k^2$ in the $\theta$-fiber spectrum of the
perfect guide.  The Floquet direct-integral decomposition would then give
$k^2\in\sigma(A^\pm(\beta))\subset\Sigma_{\mathrm{ess}}(\beta)$, contrary to
\eqref{eq:strict-gap}.  This is the spectral meaning of the absence of
propagating half-guide channels in a common essential gap.
\end{proof}

A generalized Bloch trace chain for a multiplier $\lambda\in\Complex$ is a finite
sequence $e_0,\ldots,e_{m-1}\in\mathscr T^\pm$ satisfying
\begin{equation}
  (T^\pm-\lambda I)e_0=0,
  \qquad (T^\pm-\lambda I)e_\ell=e_{\ell-1},
  \quad 1\leq\ell<m.
  \label{eq:Jordan-chain}
\end{equation}
The corresponding fields are generalized Bloch modes.  If
$\sigma(T^\pm)\cap\Sone=\varnothing$, choose a contour $\Gamma_s$ enclosing
precisely the multipliers with $|\lambda|<1$ and define
\begin{equation}
  P_s^\pm:=\frac{1}{2\pi\ii}\int_{\Gamma_s}(zI-T^\pm)^{-1}\dd z,
  \qquad \mathscr E_s^\pm:=\ran P_s^\pm.
  \label{eq:stable-projection}
\end{equation}
The subspace $\mathscr E_s^\pm$ is the stable spectral subspace.

Under a Riesz-basis hypothesis, collect all stable chains in a coefficient
Hilbert space
\begin{equation}
  \mathcal K_s^\pm:=
  \ell^2\bigl(\{(j,\ell): |\lambda_j^\pm|<1,
  0\leq\ell<m_j^\pm\}\bigr).
  \label{eq:stable-coefficients}
\end{equation}
Let
\begin{equation}
  \mathcal E^\pm(k,\beta):\mathcal K_s^\pm
  \longrightarrow\mathcal U^\pm_{\mathrm{out}}(k,\beta)
  \label{eq:field-synthesis}
\end{equation}
be the Riesz synthesis of the associated generalized fields, and define the
complete Cauchy-trace synthesis
\begin{equation}
  Q^\pm(k,\beta):=\Tr^{0,\pm}\mathcal E^\pm(k,\beta):
  \mathcal K_s^\pm\longrightarrow\mathcal H_\Gamma^\pm.
  \label{eq:trace-synthesis}
\end{equation}

\begin{proposition}[Cauchy relation, DtN graph, and stable traces]
\label{prop:relation-Bloch}
Assume the periodic half-guide and trace problem satisfy the hypotheses of the
generalized Floquet Riesz-basis theorem in
\cite[Theorem~2.2]{HohageSoussi2013}, and assume
$\sigma(T^\pm)\cap\Sone=\varnothing$.  Then
\begin{equation}
  \mathcal C^\pm_{\mathrm{out}}(k,\beta)=\ran Q^\pm(k,\beta),
  \label{eq:relation-stable-range}
\end{equation}
where the range is closed and the synthesis includes every stable Jordan
chain.  If in addition $\pi_D^\pm|_{\Coutpm}$ is bijective, then
\begin{equation}
  \mathcal C^\pm_{\mathrm{out}}(k,\beta)
  =\graph\Lambda^\pm(k,\beta).
  \label{eq:relation-DtN-graph}
\end{equation}
Ordinary Bloch eigenvectors may replace the chains only when the stable part
of $T^\pm$ is semisimple.
\end{proposition}
\status{known modal result with an adaptation required for the present transmission traces.}
\begin{proofroadmap}
First identify a unit-modulus multiplier with a real transverse Bloch
quasimomentum.  The common essential gap excludes such a perfect-guide Bloch
solution and provides separated stable and unstable Riesz contours.  Apply
\cite[Theorem~2.2 and Propositions~6.1--6.2]{HohageSoussi2013} to obtain a
Riesz basis of generalized modes and, when the trace problem is well posed, a
Riesz basis of their traces.  The spectral-decomposition route in
\cite[Sections~3--5]{Zhang2021} and the modal DtN construction in
\cite{Zhang2023} provide complementary formulations.

The nontrivial adaptation is to verify the coefficient, interface, and trace
isomorphism hypotheses for the piecewise-constant transmission half-guide used
here, including the two-component center-oriented Cauchy trace.  The graph
identity then follows immediately from \eqref{eq:DtN-definition}.  Consult
\path{research/planning/muller-cauchy/notes/bloch-trans.md}
and \path{research/planning/muller-cauchy/notes/outgoing.md}.
If a trace Riesz basis cannot be obtained, the weakest acceptable result is
the closed-range identity \eqref{eq:relation-stable-range} with coefficients
factored by $\ker Q^\pm$.
\end{proofroadmap}
\begin{proof}
It suffices to argue for one sign, denoted by
$\star\in\{-,+\}$, because the outward coordinate
$t^\star=s^\star(x-X^\star)$ puts both half-guides into the same right-going
half-strip form.

We first match the hypotheses of
\cite[Theorem~2.2 and Proposition~3.5]{HohageSoussi2013} in the order in
which they enter the argument.
\begin{enumerate}[label=\textup{(\arabic*)},leftmargin=*]
  \item Their scalar half-strip equation requires a positive, bounded,
  periodic coefficient.  After rescaling $t^\star$ by $L^\star$, this is
  precisely the periodic extension of \eqref{eq:material}; piecewise-constant
  transmission coefficients are allowed because the cited theorem assumes
  only an $L^\infty$ coefficient bounded away from zero.
  \item Their lateral boundary operator may be quasiperiodic.  The present
  $y$-quasiperiodicity is that case.  Any additional symmetry required by the
  cited theorem at its exceptional quasiperiodic parameters is included in
  the present proposition's explicit assumption that all of that theorem's
  hypotheses hold.
  \item Their outgoing solution space must be invariant under one-cell
  translation.  Here this is part of the translation-framework hypotheses.
  Moreover, $\sigma(T^\star)\cap\Sone=\varnothing$ removes all real-
  quasimomentum Floquet modes, so their radiating space
  $H^{1,+}$ reduces to the square-integrable $H^1$ solution space
  $\mathcal U^\star_{\mathrm{out}}(k,\beta)$.
  \item For the trace conclusion in part 2 of their theorem, the boundary
  problem for a trace
  $\tau=\theta_D\gamma_D+\theta_N\partial_{t^\star}$ must be well posed.
  This is exactly the ``trace problem'' part of the hypotheses assumed in the
  present proposition.  With the center-oriented convention,
  $\partial_{t^\star}$ is the second component of
  \eqref{eq:halfguide-trace}, so $\tau$ is a bounded projection of the present
  two-component Cauchy trace.
\end{enumerate}
Thus the solution-space and trace-space conclusions of the cited theorem,
and the Jordan-chain identification of its Proposition~3.5, apply without a
change of multiplier or orientation.

\emph{Step 1: identify the stable synthesis with all outgoing fields.}
The cited Riesz-basis theorem gives constants $a,b>0$ and a synthesis
isomorphism
\[
  \mathcal E^\star:\mathcal K_s^\star
  \longrightarrow\mathcal U^\star_{\mathrm{out}}(k,\beta)
\]
formed from the canonical generalized Floquet chains.  The unit-circle
exclusion implies that every mode in this $H^1$ solution space has a stable
multiplier and that every stable chain occurs in the synthesis.  Consequently,
using \eqref{eq:trace-synthesis} and then
\eqref{eq:outgoing-relation},
\[
  \ran Q^\star
  =\Tr^{0,\star}\ran\mathcal E^\star
  =\Tr^{0,\star}\mathcal U^\star_{\mathrm{out}}(k,\beta)
  =\mathcal C^\star_{\mathrm{out}}(k,\beta).
\]

\emph{Step 2: prove that the range is closed.}
Let $\mathcal Y_\tau^\star$ be the boundary space in which the well-posed
trace problem takes its data, and let
\[
  L_\tau^\star:\mathcal H_\Gamma^\star\to
  \mathcal Y_\tau^\star,
  \qquad L_\tau^\star(f,g)=\theta_Df+\theta_Ng.
\]
Part 2 of \cite[Theorem~2.2]{HohageSoussi2013} says that
$L_\tau^\star Q^\star$ is a Riesz synthesis isomorphism.  Hence, for some
$c_\tau>0$ and every $c\in\mathcal K_s^\star$,
\[
  c_\tau\lVert c\rVert_{\mathcal K_s^\star}
  \leq\lVert L_\tau^\star Q^\star c\rVert_{\mathcal Y_\tau^\star}
  \leq\lVert L_\tau^\star\rVert
  \lVert Q^\star c\rVert_{\mathcal H_\Gamma^\star}.
\]
The trace theorem and the upper Riesz estimate for $\mathcal E^\star$ also
give
\[
  \lVert Q^\star c\rVert_{\mathcal H_\Gamma^\star}
  \leq C\lVert c\rVert_{\mathcal K_s^\star}.
\]
Thus $Q^\star$ is bounded below as well as bounded.  In particular it is
injective and has closed range.  This proves all assertions in
\eqref{eq:relation-stable-range}.

\emph{Step 3: pass to the Dirichlet chart.}
Assume now that
$\pi_D^\star|_{\mathcal C^\star_{\mathrm{out}}}$ is bijective.  For every
$(f,g)\in\mathcal C^\star_{\mathrm{out}}$, definition
\eqref{eq:DtN-definition} gives
\[
  g=\pi_N^\star(f,g)
   =\pi_N^\star
    \bigl(\pi_D^\star|_{\mathcal C^\star_{\mathrm{out}}}\bigr)^{-1}f
   =\Lambda^\star f.
\]
Conversely, the inverse Dirichlet projection puts
$(f,\Lambda^\star f)$ in the relation for every
$f\in H^{1/2}_\beta(\Gamma^\star)$.  Hence
$\mathcal C^\star_{\mathrm{out}}=\graph\Lambda^\star$, which is
\eqref{eq:relation-DtN-graph}.

Finally, Proposition~3.5 of the cited paper identifies each root subspace with
finite Jordan chains for one-cell translation.  If a stable Jordan block has
size greater than one, its associated vectors are not eigenvectors and cannot
belong to the span of ordinary eigenvectors of that block.  Ordinary Bloch
eigenvectors can therefore replace all generalized chain vectors exactly when
the stable restriction of $T^\star$ is semisimple.
\end{proof}

% Full title: M\"uller Representation and the Coupled Operator
\section{M\"uller Representation and the Coupled Operator}
\label{sec:muller-coupling}

\subsection{Rayleigh and layer-potential conventions}

Let $\nu$ denote the unit normal on $\Upsilon^0=\partial\Omega^0$ pointing from
$\Omega^0$ into the background.  Interior and exterior traces are
\[
  \gamma^{\mathrm i}u(z):=\lim_{h\downarrow0}u(z-h\nu(z)),
  \qquad
  \gamma^{\mathrm e}u(z):=\lim_{h\downarrow0}u(z+h\nu(z)).
\]
The same superscripts are used for normal derivatives in the direction $\nu$.
The outgoing two-dimensional fundamental solution is
\begin{equation}
  \Phi_\kappa(x,z):=\frac{\ii}{4}H_0^{(1)}(\kappa|x-z|).
  \label{eq:fundamental-solution}
\end{equation}
For $x\notin\Upsilon^0$, define
\begin{align}
  (S_\kappa\sigma)(x)&:=\int_{\Upsilon^0}\Phi_\kappa(x,z)\sigma(z)\dd s(z),
  \label{eq:single-layer}\\
  (D_\kappa\tau)(x)&:=\int_{\Upsilon^0}
  \partial_{\nu(z)}\Phi_\kappa(x,z)\tau(z)\dd s(z).
  \label{eq:double-layer}
\end{align}
With these conventions,
$\gamma_D^{\mathrm e}D_\kappa-\gamma_D^{\mathrm i}D_\kappa=I$ and
$\gamma_N^{\mathrm e}S_\kappa-\gamma_N^{\mathrm i}S_\kappa=-I$.

Let $G^{\mathrm{qp}}_{k,\beta}$ be the vertically $\beta$-quasiperiodic
outgoing Green function satisfying
\[
  G^{\mathrm{qp}}_{k,\beta}(x+de_y,z)
  =e^{\ii\beta d}G^{\mathrm{qp}}_{k,\beta}(x,z),
  \qquad (\Delta_x+k^2)G^{\mathrm{qp}}_{k,\beta}(x,z)=-\delta_z.
\]
Away from its exceptional set it is represented by the lattice sum
\begin{equation}
  G^{\mathrm{qp}}_{k,\beta}(x,z)
  =\sum_{\ell\in\Z}e^{\ii\beta\ell d}
  \Phi_k(x,z+\ell d e_y),
  \label{eq:qp-Green}
\end{equation}
interpreted by the standard convergent continuation when necessary.  Replacing
$\Phi_k$ by $G^{\mathrm{qp}}_{k,\beta}$ in
\eqref{eq:single-layer}--\eqref{eq:double-layer} defines
$S^{\mathrm{qp}}_{k,\beta}$ and $D^{\mathrm{qp}}_{k,\beta}$.

For
\[
  \gamma_m(k,\beta):=\sqrt{k^2-\beta_m^2},
  \qquad \operatorname{Im}\gamma_m\geq0,
\]
a Wood anomaly occurs when $\gamma_m=0$ for some $m$.  At a non-Wood
parameter, define
\begin{equation}
  h^{1/2}_\beta:=
  \left\{a=(a_m)_{m\in\Z}:
  \sum_m\langle\beta_m\rangle|a_m|^2<\infty\right\},
  \qquad
  \Xi_\beta:=h^{1/2}_\beta\times h^{1/2}_\beta.
  \label{eq:Rayleigh-space}
\end{equation}
For $\xi=(\xi^\rightarrow,\xi^\leftarrow)\in\Xi_\beta$, the homogeneous
Rayleigh reconstruction in the center cell is
\begin{equation}
  (\mathcal H_{k,\beta}\xi)(x,y)
  :=\sum_{m\in\Z}\left[
  \xi_m^\rightarrow e^{\ii\gamma_m(x-X^-)}
  +\xi_m^\leftarrow e^{-\ii\gamma_m(x-X^+)}\right]\psi_m(y).
  \label{eq:Rayleigh-reconstruction}
\end{equation}
The weights in \eqref{eq:Rayleigh-space} control both the
$H^{1/2}_\beta$ Dirichlet trace and the $H^{-1/2}_\beta$ normal trace.

The interface-density space is
\begin{equation}
  \mathcal D^0:=H^{1/2}(\Upsilon^0)\times H^{-1/2}(\Upsilon^0),
  \qquad \eta=(\tau,\sigma)\in\mathcal D^0,
  \label{eq:density-space}
\end{equation}
where $\tau$ is the double-layer density and $\sigma$ is the single-layer
density.  To describe a field before its two material traces have been
matched, set
\begin{equation}
  H^1_{\beta,\mathrm{br}}(\mathcal C^0):=
  H^1_\beta(\mathcal C^0\setminus\overline{\Omega^0})
  \oplus H^1(\Omega^0).
  \label{eq:center-broken-space}
\end{equation}
The subscript ``br'' means \emph{broken across $\Upsilon^0$}: the two components
have independent traces until the transmission residual below is set to zero.
Let $k^0:=n^0k$.  The center reconstruction
\begin{equation}
  \mathcal R_{\mathrm c}(k,\beta):\mathcal D^0\times\Xi_\beta
  \longrightarrow H^1_{\beta,\mathrm{br}}(\mathcal C^0)
  \label{eq:center-reconstruction-map}
\end{equation}
is the broken field $u_{\eta,\xi}^{\mathrm c}$ with
\begin{equation}
  u_{\eta,\xi}^{\mathrm c}(x)=
  \begin{cases}
  \mathcal H_{k,\beta}\xi(x)
  +D^{\mathrm{qp}}_{k,\beta}\tau(x)
  -S^{\mathrm{qp}}_{k,\beta}\sigma(x),
  &x\in\mathcal C^0\setminus\overline{\Omega^0},\\
  D_{k^0}\tau(x)-S_{k^0}\sigma(x),
  &x\in\Omega^0.
  \end{cases}
  \label{eq:center-reconstruction}
\end{equation}
Thus the exterior component is vertically quasiperiodic, while the bounded
interior component needs no separate quasiperiodic condition.

Define the M\"uller transmission residual by taking all traces in the fixed
normal direction $\nu$:
\begin{equation}
\begin{split}
  \mathcal M(k,\beta)(\eta,\xi):=
  \bigl(&\gamma_D^{\mathrm e}u_{\eta,\xi}^{\mathrm c}
         -\gamma_D^{\mathrm i}u_{\eta,\xi}^{\mathrm c},\\
        &\gamma_N^{\mathrm e}u_{\eta,\xi}^{\mathrm c}
         -\gamma_N^{\mathrm i}u_{\eta,\xi}^{\mathrm c}\bigr)
  \in\mathcal H_{\Upsilon^0},
  \qquad
  \mathcal H_{\Upsilon^0}:=H^{1/2}(\Upsilon^0)\times H^{-1/2}(\Upsilon^0).
\end{split}
\label{eq:Muller-residual}
\end{equation}
This trace definition fixes every half-identity sign without relying on an
unwritten matrix convention.  Its expansion using the jump relations is the
M\"uller block.  Let
\begin{equation}
  \mathcal Z_{\mathrm c}(k,\beta):=\ker\mathcal M(k,\beta),
  \qquad
  \mathcal N_{\mathrm c}(k,\beta):=
  \mathcal Z_{\mathrm c}(k,\beta)\cap\ker\mathcal R_{\mathrm c}(k,\beta).
  \label{eq:center-nullspaces}
\end{equation}
Thus $\mathcal N_{\mathrm c}$ contains exactly the center algebraic
representations that satisfy the transmission equations but generate the zero
physical center field.

Let $\mathcal S_{\mathrm c}(k,\beta)$ be the space of piecewise Helmholtz
fields in $H^1_\beta(\mathcal C^0)$ satisfying the physical TM transmission
conditions at $\Upsilon^0$, with no boundary condition at $\Gamma^\pm$.
Define the representation exceptional set $\mathfrak E_{\mathrm{rep}}(\beta)$
as the union of
\begin{enumerate}[label=(\roman*),leftmargin=2.2em]
  \item Wood parameters $\gamma_m(k,\beta)=0$;
  \item poles or undefined parameters of the chosen quasiperiodic Green
  function;
  \item parameters for which the complementary swapped-wavenumber transmission
  problem with interior wavenumber $k$ and exterior wavenumber $k^0$ is not
  uniquely solvable.
\end{enumerate}
The quotient below retains possible representation nullspaces even outside
this explicitly listed set.

\begin{theorem}[Center-cell M\"uller--Rayleigh representation]
\label{thm:center-representation}
Let $k^2\in I_\beta$ and $k\notin\mathfrak E_{\mathrm{rep}}(\beta)$.  The
restriction of $\mathcal R_{\mathrm c}(k,\beta)$ to
$\mathcal Z_{\mathrm c}(k,\beta)$ is surjective onto
$\mathcal S_{\mathrm c}(k,\beta)$.  Consequently it induces a linear
isomorphism
\begin{equation}
  \mathcal Z_{\mathrm c}(k,\beta)/\mathcal N_{\mathrm c}(k,\beta)
  \xrightarrow{\ \cong\ }\mathcal S_{\mathrm c}(k,\beta),
  \qquad [\eta,\xi]\longmapsto\mathcal R_{\mathrm c}(k,\beta)(\eta,\xi).
  \label{eq:center-quotient-isomorphism}
\end{equation}
No uniqueness of the densities is asserted.  Density uniqueness is equivalent
to the additional statement $\mathcal N_{\mathrm c}(k,\beta)=\{0\}$.
\end{theorem}
\status{adaptation and proof target.}
\begin{proofroadmap}
Extract from a given center field the two homogeneous Rayleigh components in
homogeneous collars of $\Gamma^-$ and $\Gamma^+$.  Subtract their continuation
\eqref{eq:Rayleigh-reconstruction}; the remainder satisfies the outgoing
single-cell representation hypotheses.  Apply the quasiperiodic exterior and
free-space interior Green identities and construct the complementary
swapped-wavenumber field as in
\cite[Theorem~4 and Appendix~A, Lemmas~9--10]{BarnettGreengard2010}.  The
exterior double layer in this calculation is
$D^{\mathrm{qp}}_{k,\beta}$, whereas the inclusion double layer is $D_{k^0}$;
using $D^{\mathrm{qp}}_{k^0,\beta}$ in the exterior extinction identity is the
kernel mismatch that must not be repeated.

Surjectivity of the physical-field map is the substantive claim.  Once it is
proved, the quotient isomorphism follows from the first isomorphism theorem.
The density kernel must be analyzed separately with jump relations,
complementary fields, and Calder\'on projector ranges; the transmission-BIE
analysis in \cite[Lemmas~2.2--2.6]{HiptmairMoiolaSpence2022} explains why a
physical uniqueness theorem does not automatically prove density uniqueness.

Detailed internal guidance is in
\path{research/planning/muller-cauchy/notes/layers.md} and the
obligations named ``corrected center representation'' and
``density-to-field kernel characterization'' in
\path{research/planning/muller-cauchy/p-proofs.md}.  The weakest
acceptable result is surjectivity modulo a finite-dimensional missed
homogeneous sector; an unquotiented theorem is acceptable only after
$\mathcal N_{\mathrm c}=\{0\}$ is proved on the stated regular set.
\end{proofroadmap}
\begin{proof}
Write $u=(u^{\mathrm e},u^{\mathrm i})\in
\mathcal S_{\mathrm c}(k,\beta)$ for the exterior and inclusion components,
and construct $(\eta,\xi)\in\mathcal Z_{\mathrm c}(k,\beta)$ whose
reconstructed field is $u$.

\emph{Step 1: separate the incoming homogeneous Rayleigh field.}
Because $\Omega^0\Subset\mathcal C^0$, there are homogeneous collars of both
ports.  In either collar, the Fourier coefficient of $u^{\mathrm e}$ in the
mode $\psi_m$ solves
\[
  \partial_x^2u_m+\gamma_m^2u_m=0.
\]
The non-Wood assumption gives $\gamma_m\neq0$, so the solution has a unique
right/left-going decomposition.  Denote by $a_m^{-,\rightarrow}$ the
right-going coefficient in the left collar and by
$a_m^{+,\leftarrow}$ the left-going coefficient in the right collar, with
the phases based respectively at $X^-$ and $X^+$.  If $(f_m^\pm,p_m^\pm)$
are the Fourier coefficients of the Dirichlet and $\partial_x$ traces at the
two ports, then
\[
  a_m^{-,\rightarrow}
  =\frac12\left(f_m^-+\frac{p_m^-}{\ii\gamma_m}\right),
  \qquad
  a_m^{+,\leftarrow}
  =\frac12\left(f_m^+-\frac{p_m^+}{\ii\gamma_m}\right).
\]
The trace theorem gives $f^\pm\in h^{1/2}_\beta$ and
$p^\pm\in h^{-1/2}_\beta$.  Moreover,
$|\gamma_m|\asymp\langle\beta_m\rangle$ for all sufficiently large $|m|$;
the remaining set is finite and contains no zero $\gamma_m$.  Hence
\[
  \xi:=\bigl((a_m^{-,\rightarrow})_m,
             (a_m^{+,\leftarrow})_m\bigr)\in\Xi_\beta.
\]
Set $h:=\mathcal H_{k,\beta}\xi$ and $w:=u^{\mathrm e}-h$.  In the left
collar $w$ contains only left-going modes, and in the right collar it contains
only right-going modes: the term launched from the opposite port contributes
only to the corresponding outgoing component.  Continuing these modal
series through the two ports therefore extends $w$ to a vertically
$\beta$-quasiperiodic solution of $(\Delta+k^2)w=0$ in
$B\setminus\overline{\Omega^0}$ which is outgoing at both ends.

\emph{Step 2: the two Green representation identities.}
Use the fixed normal $\nu$ and abbreviate
\[
  P_\kappa(f,q):=D_\kappa f-S_\kappa q,
  \qquad
  P^{\mathrm{qp}}_{k,\beta}(f,q)
  :=D^{\mathrm{qp}}_{k,\beta}f-S^{\mathrm{qp}}_{k,\beta}q.
\]
Green's second identity, first on a truncated quasiperiodic strip and then by
the outgoing limiting-absorption limit, gives
\[
  P^{\mathrm{qp}}_{k,\beta}
  (\gamma_D^{\mathrm e}w,\gamma_N^{\mathrm e}w)
  =w
  \quad\text{in }B\setminus\overline{\Omega^0}.
\]
The same identity applied to a solution $z$ of $(\Delta+k^2)z=0$ in
$\Omega^0$ yields the extinction relation
\[
  P^{\mathrm{qp}}_{k,\beta}
  (\gamma_D^{\mathrm i}z,\gamma_N^{\mathrm i}z)=0
  \quad\text{in }B\setminus\overline{\Omega^0}.
\]
For the free-space kernel at $k^0$, the corresponding identities are
\begin{align*}
  P_{k^0}(\gamma_D^{\mathrm i}z,\gamma_N^{\mathrm i}z)
  &=-z&&\text{in }\Omega^0,\\
  P_{k^0}(\gamma_D^{\mathrm e}z,\gamma_N^{\mathrm e}z)
  &=0&&\text{in }\Omega^0
\end{align*}
when the latter $z$ is an outgoing exterior field.  The signs follow from
the jump conventions preceding \eqref{eq:single-layer}.  These are the
one-periodic analogues of the identities used in
\cite[Appendix~A, Lemmas~9--10]{BarnettGreengard2010}; here the vertical
wall terms cancel by $\beta$-quasiperiodicity and the two end terms are
removed by outgoing limiting absorption.

\emph{Step 3: construct common M\"uller densities.}
Let
\[
  f:=\gamma_D^{\mathrm e}u^{\mathrm e}
    =\gamma_D^{\mathrm i}u^{\mathrm i},
  \qquad
  q:=\gamma_N^{\mathrm e}u^{\mathrm e}
    =\gamma_N^{\mathrm i}u^{\mathrm i},
\]
and write $h_D:=\gamma_Dh$, $h_N:=\gamma_Nh$ on $\Upsilon^0$.  Then the
exterior Cauchy data of $w$ are $(f-h_D,q-h_N)$.  Consider the complementary
swapped-wavenumber transmission problem
\begin{align*}
  (\Delta+k^2)v^{\mathrm i}&=0
    &&\text{in }\Omega^0,\\
  (\Delta+(k^0)^2)v^{\mathrm e}&=0
    &&\text{in }\R^2\setminus\overline{\Omega^0},
\end{align*}
with the Sommerfeld condition for $v^{\mathrm e}$ and jumps
\begin{align*}
  \gamma_D^{\mathrm e}v^{\mathrm e}
  -\gamma_D^{\mathrm i}v^{\mathrm i}&=2f-h_D,\\
  \gamma_N^{\mathrm e}v^{\mathrm e}
  -\gamma_N^{\mathrm i}v^{\mathrm i}&=2q-h_N.
\end{align*}
Clause (iii) in the definition of
$\mathfrak E_{\mathrm{rep}}(\beta)$ and
$k\notin\mathfrak E_{\mathrm{rep}}(\beta)$ give its unique solution for
these admissible trace data.  Define
\begin{align*}
  \tau
  &:=f-h_D+\gamma_D^{\mathrm i}v^{\mathrm i}
    =-f+\gamma_D^{\mathrm e}v^{\mathrm e},\\
  \sigma
  &:=q-h_N+\gamma_N^{\mathrm i}v^{\mathrm i}
    =-q+\gamma_N^{\mathrm e}v^{\mathrm e}.
\end{align*}
The equalities are exactly the two prescribed jumps, and the trace mapping
properties put $(\tau,\sigma)$ in $\mathcal D^0$.

Set $\eta=(\tau,\sigma)$.  By the two quasiperiodic identities above, the
exterior component of the reconstruction is
\begin{align*}
  h+P^{\mathrm{qp}}_{k,\beta}(\tau,\sigma)
  &=h+P^{\mathrm{qp}}_{k,\beta}(f-h_D,q-h_N)\\
  &\quad
    +P^{\mathrm{qp}}_{k,\beta}
      (\gamma_D^{\mathrm i}v^{\mathrm i},
       \gamma_N^{\mathrm i}v^{\mathrm i})\\
  &=h+w=u^{\mathrm e}.
\end{align*}
Likewise, using the alternative expressions for $\tau,\sigma$ and the two
free-space identities above, its interior component is
\begin{align*}
  P_{k^0}(\tau,\sigma)
  &=-P_{k^0}(f,q)
    +P_{k^0}(\gamma_D^{\mathrm e}v^{\mathrm e},
             \gamma_N^{\mathrm e}v^{\mathrm e})\\
  &=u^{\mathrm i}.
\end{align*}
Thus $\mathcal R_{\mathrm c}(k,\beta)(\eta,\xi)=u$.  Since $u$ satisfies
the physical transmission conditions, its two jumps vanish, so
$\mathcal M(k,\beta)(\eta,\xi)=0$ and
$(\eta,\xi)\in\mathcal Z_{\mathrm c}(k,\beta)$.  This proves the asserted
surjectivity.

\emph{Step 4: pass to the quotient.}
The kernel of
$\mathcal R_{\mathrm c}|_{\mathcal Z_{\mathrm c}}$ is, by
\eqref{eq:center-nullspaces}, exactly $\mathcal N_{\mathrm c}$.  The first
isomorphism theorem therefore gives
\eqref{eq:center-quotient-isomorphism}.  This argument neither requires nor
asserts that $\mathcal N_{\mathrm c}$ is zero.
\end{proof}

\subsection{The continuous coupled operator}

The port trace of the center reconstruction is the bounded map
\begin{equation}
  \Pi^\pm(k,\beta):=\Tr^{0,\pm}\mathcal R_{\mathrm c}(k,\beta):
  \mathcal D^0\times\Xi_\beta\longrightarrow\mathcal H_\Gamma^\pm.
  \label{eq:center-port-map}
\end{equation}
This definition includes both the layer contribution and the homogeneous
Rayleigh contribution.  The half-guide maps $\mathcal E^\pm$ and $Q^\pm$ were
defined in \eqref{eq:field-synthesis}--\eqref{eq:trace-synthesis}.

\begin{definition}[Coupled M\"uller--Cauchy operator and representation nullspace]
\label{def:coupled-operator}
Define the algebraic domain and residual codomain by
\begin{align}
  \mathcal X(k,\beta)
  &:=\mathcal D^0\times\Xi_\beta\times
  \mathcal K_s^-\times\mathcal K_s^+,
  \label{eq:coupled-domain}\\
  \mathcal Y(k,\beta)
  &:=\mathcal H_{\Upsilon^0}\times
  \mathcal H_\Gamma^-\times\mathcal H_\Gamma^+.
  \label{eq:coupled-codomain}
\end{align}
For $x=(\eta,\xi,c^-,c^+)\in\mathcal X(k,\beta)$, set
\begin{equation}
  \mathbb A(k,\beta)x:=
  \begin{bmatrix}
  \mathcal M(k,\beta)(\eta,\xi)\\
  \Pi^-(k,\beta)(\eta,\xi)-Q^-(k,\beta)c^-\\
  \Pi^+(k,\beta)(\eta,\xi)-Q^+(k,\beta)c^+
  \end{bmatrix}
  \in\mathcal Y(k,\beta).
  \label{eq:coupled-operator}
\end{equation}
The first row is the physical transmission mismatch on $\Upsilon^0$; the second
and third rows are the complete center-oriented Cauchy mismatches at the left
and right ports.

Let
\begin{equation}
  H^1_{\beta,\mathrm{br}}(B):=
  H^1_\beta(W^- )\oplus
  H^1_\beta(\mathcal C^0\setminus\overline{\Omega^0})
  \oplus H^1(\Omega^0)
  \oplus H^1_\beta(W^+).
  \label{eq:global-broken-space}
\end{equation}
The global broken-field reconstruction is
\begin{equation}
  \mathcal R_{\mathrm g}(k,\beta)x:=
  \begin{cases}
    \mathcal E^-(k,\beta)c^-,&\text{in }W^-,\\
    \mathcal R_{\mathrm c}(k,\beta)(\eta,\xi),&\text{in }\mathcal C^0,\\
    \mathcal E^+(k,\beta)c^+,&\text{in }W^+.
  \end{cases}
  \label{eq:global-reconstruction}
\end{equation}
It is a broken field for general $x$ and a global $H^1_\beta(B)$ field when
$x\in\ker\mathbb A(k,\beta)$.  Finally, the representation-induced nullspace
is
\begin{equation}
  \mathcal N_{\mathrm{rep}}(k,\beta):=
  \left\{x\in\ker\mathbb A(k,\beta):
  \mathcal R_{\mathrm g}(k,\beta)x=0
  \text{ in }H^1_{\beta,\mathrm{br}}(B)\right\}.
  \label{eq:Nrep-definition}
\end{equation}
Thus $\mathbb A$ is the fully defined algebraic coupling, while
$\mathcal N_{\mathrm{rep}}$ records precisely its kernel vectors that carry no
physical field.  Neither symbol is shorthand for an unspecified numerical
matrix.
\end{definition}

% Full title: Main Results and Proof Architecture
\section{Main Results and Proof Architecture}
\label{sec:main-results}

We collect the assumptions required by the target theorem.  They define the
\emph{regular parameter set} $\mathfrak R$ of pairs $(k,\beta)$ for which:
\begin{enumerate}[label=(R\arabic*),leftmargin=2.8em]
  \item $k>0$ is real and $k^2$ lies in the compact gap window
  $I_\beta$ of \eqref{eq:strict-gap};
  \item $k\notin\mathfrak E_{\mathrm{rep}}(\beta)$ and all layer, Rayleigh,
  and port-trace maps in \Cref{sec:muller-coupling} are bounded on the stated
  spaces;
  \item the half-guide translations have no unit-circle spectrum and
  \Cref{prop:relation-Bloch} holds with injective Riesz synthesis maps
  $\mathcal E^\pm$;
  \item the center representation in \Cref{thm:center-representation} is
  surjective and $\mathcal N_{\mathrm c}(k,\beta)$ is closed.
\end{enumerate}
These assumptions deliberately exclude thresholds, Wood anomalies, embedded
guided states, complex resonances, and representation poles.

Before stating the main equation, it is useful to keep its three objects at
different conceptual levels:
\begin{enumerate}[label=(\alph*),leftmargin=2.2em]
  \item $\mathbb A(k,\beta):\mathcal X(k,\beta)\to\mathcal Y(k,\beta)$ is the
  \emph{residual operator} defined in \eqref{eq:coupled-operator}.  Its input
  is one tuple $x=(\eta,\xi,c^-,c^+)$ of center interface densities,
  homogeneous Rayleigh coefficients, and left/right stable-mode
  coefficients.  The equation $\mathbb A x=0$ says exactly that the center
  transmission jumps and both center--half-guide Cauchy mismatches vanish.
  \item $\mathcal N_{\mathrm{rep}}(k,\beta)$ is the subset of
  $\ker\mathbb A(k,\beta)$ defined in \eqref{eq:Nrep-definition} whose
  densities and modal coefficients reconstruct the zero field everywhere.
  It is removed because two algebraic descriptions that generate the same
  physical field must represent the same solution class.
  \item $\mathcal G(k,\beta)=\ker(A(\beta)-k^2I)$ is the physical guided-mode
  eigenspace defined in \eqref{eq:guided-space}: its elements are global,
  vertically $\beta$-quasiperiodic $H^1$ solutions on the infinite strip and
  are therefore square-integrable in the transverse direction.
\end{enumerate}
Thus the theorem does not identify all algebraic vectors with fields one by
one; it identifies algebraic kernel vectors modulo their zero-field
representations with the physical eigenspace.

\begin{theorem}[Kernel--guided-mode equivalence: main target]
\label{thm:main-kernel-equivalence}
Let $(k,\beta)\in\mathfrak R$.  The reconstruction
$\mathcal R_{\mathrm g}(k,\beta)$ maps
$\ker\mathbb A(k,\beta)$ onto the global guided-mode eigenspace
$\mathcal G(k,\beta)$, and
\begin{equation}
  \ker\!\left(
  \mathcal R_{\mathrm g}(k,\beta)|_{\ker\mathbb A(k,\beta)}
  \right)
  =\mathcal N_{\mathrm{rep}}(k,\beta)
  =\mathcal N_{\mathrm c}(k,\beta)\times\{0\}\times\{0\}.
  \label{eq:representation-kernel-identification}
\end{equation}
Consequently, reconstruction induces a natural linear isomorphism
\begin{equation}
  \overline{\mathcal R}_{\mathrm g}(k,\beta):
  \frac{\ker\mathbb A(k,\beta)}{\mathcal N_{\mathrm{rep}}(k,\beta)}
  \xrightarrow{\ \cong\ }\mathcal G(k,\beta).
  \label{eq:main-quotient-isomorphism}
\end{equation}
In particular,
\begin{equation}
  \dim\!\left(
  \ker\mathbb A(k,\beta)/\mathcal N_{\mathrm{rep}}(k,\beta)
  \right)=m_{\mathrm g}(k,\beta).
  \label{eq:multiplicity-preservation}
\end{equation}
If the stronger density-injectivity statement
$\mathcal N_{\mathrm c}(k,\beta)=\{0\}$ holds, then the quotient can be
removed and $\ker\mathbb A(k,\beta)\cong\mathcal G(k,\beta)$.
\end{theorem}
\status{main proof target.}
\begin{proofroadmap}
For completeness, take $u\in\mathcal G(k,\beta)$.  By
\Cref{thm:bounded-reduction}, its center restriction belongs to
$\mathcal S_{\mathrm c}(k,\beta)$ and its half-guide restrictions lie in the two
outgoing spaces.  Apply \Cref{thm:center-representation} to choose a center
class $[\eta,\xi]$, and apply \eqref{eq:relation-stable-range} to obtain the
unique generalized-mode coordinates $c^\pm$.  The physical interface and port
matching conditions then give $x=(\eta,\xi,c^-,c^+)\in\ker\mathbb A$.

For soundness, start with $x\in\ker\mathbb A$.  The first row of
\eqref{eq:coupled-operator} makes the center reconstruction a physical
transmission field.  The other rows identify its complete Cauchy data with
those of the two synthesized half-guide fields.  Apply
\Cref{lem:weak-gluing} and the strict-gap $H^1$ selection to obtain a member of
$\mathcal G(k,\beta)$.

For injectivity, assume the reconstructed global field is zero.  Injectivity
of the generalized Riesz synthesis gives $c^-=c^+=0$.  The center pair then
lies in $\mathcal N_{\mathrm c}$, proving
\eqref{eq:representation-kernel-identification}.  This last step is the main
proof bottleneck if the half-guide synthesis is only a closed-range map or if the
center density kernel has not been characterized.  The closest prior
ingredients are \cite[Theorem~4.5]{Fliss2013},
\cite[Theorem~4 and Appendix~A]{BarnettGreengard2010}, and
\cite[Theorem~2.2]{HohageSoussi2013}; none by itself proves the coupled
statement.

The full internal obligation is headed ``continuous kernel--guided-space
isomorphism'' in
\path{research/planning/muller-cauchy/p-proofs.md}; the dependency
analysis is in
\path{research/planning/muller-cauchy/p-deps.md}.  The
weakest acceptable conclusion is a surjection from the quotient kernel onto
$\mathcal G(k,\beta)$ together with an abstract, rather than explicit,
description of $\mathcal N_{\mathrm{rep}}$.
\end{proofroadmap}
\begin{proof}
Fix $(k,\beta)\in\mathfrak R$ and suppress these parameters from the
notation.

\emph{Step 1: reconstruction maps the algebraic kernel into the guided-mode
space.}
Let $x=(\eta,\xi,c^-,c^+)\in\ker\mathbb A$.  Its first block equation is
$\mathcal M(\eta,\xi)=0$.  Hence the two material traces of
$\mathcal R_{\mathrm c}(\eta,\xi)$ agree, and the center reconstruction is a
field in $\mathcal S_{\mathrm c}$.  The two remaining block equations say
\[
  \Pi^\pm(\eta,\xi)=Q^\pm c^\pm
  =\Tr^{0,\pm}\mathcal E^\pm c^\pm.
\]
Thus the center field and each synthesized half-guide field have identical
center-oriented Cauchy data at their common port.  Since
$\mathcal E^\pm c^\pm\in\mathcal U^\pm_{\mathrm{out}}$, two applications of
\Cref{lem:weak-gluing} show that $\mathcal R_{\mathrm g}x$ is a global field
in $H^1_\beta(B)$ satisfying the weak eigenvalue equation.  By
\Cref{prop:weak-strong},
\[
  \mathcal R_{\mathrm g}x\in\mathcal G(k,\beta).
\]
Consequently the restriction of $\mathcal R_{\mathrm g}$ to
$\ker\mathbb A$ is a well-defined linear map into $\mathcal G(k,\beta)$.

\emph{Step 2: every guided mode has algebraic coordinates.}
Take $u\in\mathcal G(k,\beta)$.  By
\Cref{thm:bounded-reduction}, its center restriction
$u^0:=u|_{\mathcal C^0}$ is a physical center solution and the restrictions
$u^\pm:=u|_{W^\pm}$ belong to
$\mathcal U^\pm_{\mathrm{out}}$.  Assumption (R4) and
\Cref{thm:center-representation} provide
$(\eta,\xi)\in\mathcal Z_{\mathrm c}$ such that
\[
  \mathcal R_{\mathrm c}(\eta,\xi)=u^0.
\]
By assumption (R3) and \eqref{eq:relation-stable-range}, there are
$c^\pm\in\mathcal K_s^\pm$ such that
\[
  Q^\pm c^\pm=\Tr^{0,\pm}u^\pm.
\]
The converse part of \Cref{lem:weak-gluing} gives
$\Tr^{0,\pm}u^0=\Tr^{0,\pm}u^\pm$.  Therefore
$x=(\eta,\xi,c^-,c^+)$ makes all three rows of
\eqref{eq:coupled-operator} vanish.  It lies in $\ker\mathbb A$, and its
piecewise reconstruction is exactly $u$.  Hence
$\mathcal R_{\mathrm g}|_{\ker\mathbb A}$ is onto
$\mathcal G(k,\beta)$.

\emph{Step 3: identify the representation kernel.}
The first equality in \eqref{eq:representation-kernel-identification} follows
immediately from the definition \eqref{eq:Nrep-definition}.  We prove the
second.  Suppose $x\in\ker\mathbb A$ and
$\mathcal R_{\mathrm g}x=0$.  On the two half-guides this says
\[
  \mathcal E^-c^-=0,
  \qquad
  \mathcal E^+c^+=0.
\]
The injectivity required in (R3) gives $c^-=c^+=0$.  On the center cell,
$\mathcal R_{\mathrm c}(\eta,\xi)=0$, while the first row of
$\mathbb A x=0$ gives $(\eta,\xi)\in\mathcal Z_{\mathrm c}$.  Thus
$(\eta,\xi)\in\mathcal N_{\mathrm c}$ and
\[
  \mathcal N_{\mathrm{rep}}
  \subset\mathcal N_{\mathrm c}\times\{0\}\times\{0\}.
\]
Conversely, if $(\eta,\xi)\in\mathcal N_{\mathrm c}$ and $c^-=c^+=0$,
then $\mathcal M(\eta,\xi)=0$ and
$\Pi^\pm(\eta,\xi)=\Tr^{0,\pm}\mathcal R_{\mathrm c}(\eta,\xi)=0$.
Hence $(\eta,\xi,0,0)\in\ker\mathbb A$ and reconstructs the zero global
field.  The reverse inclusion follows, proving
\eqref{eq:representation-kernel-identification}.

\emph{Step 4: quotient and multiplicity.}
We have proved that the linear map
$\mathcal R_{\mathrm g}|_{\ker\mathbb A}$ is surjective and has kernel
$\mathcal N_{\mathrm{rep}}$.  The first isomorphism theorem gives
\eqref{eq:main-quotient-isomorphism}.  Taking dimensions and using
\eqref{eq:geometric-multiplicity} yields
\eqref{eq:multiplicity-preservation}.  Finally, if
$\mathcal N_{\mathrm c}=\{0\}$, then
\eqref{eq:representation-kernel-identification} makes the restricted
reconstruction injective, so no quotient is needed.
\end{proof}

\subsection{Interpretation and the Fredholm question}

The term ``spurious-free'' is justified at the continuous level only after
\eqref{eq:representation-kernel-identification}: a kernel class then carries a
nonzero guided field unless it is the zero class.  This does not control large
inverse norms or discretization-induced near-kernels.  Those distinct issues
are emphasized by the spurious quasi-resonance results in
\cite[Theorems~1.10--1.12]{HiptmairMoiolaSpence2022}.

If $\mathcal N_{\mathrm c}$ is closed, let
\[
  \widehat{\mathcal X}(k,\beta):=
  \mathcal X(k,\beta)/
  \bigl(\mathcal N_{\mathrm c}(k,\beta)\times\{0\}\times\{0\}\bigr).
\]
Since this nullspace is contained in $\ker\mathbb A$, the coupled operator
descends to a bounded map
$\widehat{\mathbb A}(k,\beta):\widehat{\mathcal X}(k,\beta)
\to\mathcal Y(k,\beta)$.  Recall that a bounded operator between Banach spaces
is Fredholm if its range is closed and its kernel and cokernel are finite
dimensional; its index is $\ind F=\dim\ker F-\dim\operatorname{coker}F$.

\begin{conjecture}[Fredholm index-zero realization]
\label{conj:Fredholm}
There is a natural square multitrace or quotient realization of
$\widehat{\mathbb A}(k,\beta)$ on the regular parameter set for which
\begin{equation}
  \widehat{\mathbb A}(k,\beta)
  =\mathbb A_0+\mathbb K(k,\beta),
  \qquad \mathbb A_0^{-1}\text{ bounded},\quad
  \mathbb K(k,\beta)\text{ compact},
  \label{eq:Fredholm-splitting}
\end{equation}
and hence $\widehat{\mathbb A}(k,\beta)$ is Fredholm of index zero.
\end{conjecture}
\status{currently conjectural strengthening.}
\begin{proofroadmap}
The principal candidate for $\mathbb A_0$ combines the M\"uller/Calder\'on
principal block with identity blocks obtained by Riesz-normalizing the stable
trace syntheses.  Smooth quasiperiodic-minus-free-space kernel differences and
trace maps between separated curves are compact candidates.  Every remainder
must be checked on the precise Sobolev products in
\eqref{eq:coupled-domain}--\eqref{eq:coupled-codomain}; a block display that is
not square does not have a Fredholm index.

See \path{research/planning/muller-cauchy/notes/fredholm.md} and
the obligations named ``reference operator plus compact perturbation'' and
``Fredholm index zero'' in
\path{research/planning/muller-cauchy/p-proofs.md}.  The closest
structural precedents are \cite[Theorem~4]{BarnettGreengard2010} and
\cite[Lemma~5.1]{HohageSoussi2013}.  If no natural square realization exists,
the conjecture must be replaced by a closed-range or semi-Fredholm statement,
or by an index-zero theorem only after finite modal truncation.  The main
quotient equivalence in \Cref{thm:main-kernel-equivalence} does not logically
depend on this conjecture.
\end{proofroadmap}

\appendix
% Full title: Essential Spectrum for Two Different Periodic Ends
% Purpose:
%   Prove the essential-spectrum formula for a scalar strip with two different
%   periodic ends and supply every localization lemma used by the proof.
% Main workflow:
%   Establish local compactness and cutoff commutator bounds, derive a Zhislin
%   criterion, construct one-sided periodic Weyl sequences, and prove both
%   inclusions in the essential-spectrum identity.
% Based on:
%   The Zhislin-sequence mechanism of De Nittis et al. (2021) and the periodic
%   Floquet--Bloch characterization of Fliss (2013).
% Main changes:
%   Adapts the one-dimensional first-order argument to the present
%   two-dimensional second-order weighted strip operator with exact periodic
%   tails and possibly different transverse periods.
% Mathematical goal:
%   Prove Theorem~\ref{thm:two-ended-essential-spectrum} without assuming the
%   essential-spectrum union formula.

\section{Essential spectrum for two different periodic ends}
\label{app:essential-spectrum}

Throughout this appendix, fix
$\beta\in[-\pi/d,\pi/d)$ and abbreviate
\[
  A:=A(\beta),\qquad A^\star:=A^\star(\beta),
  \qquad \star\in\{-,+\}.
\]
The operators act respectively in
$H=L^2(B,\rho\,\dd x\dd y)$ and
$H^\star=L^2(B,\rho^\star\,\dd x\dd y)$.  Define
\begin{equation}
  H_0:=H^-\oplus H^+,
  \qquad A_0:=A^-\oplus A^+.
  \label{eq:appendix-free-operator}
\end{equation}
Let $Q$ denote multiplication by $x$ on a strip space and let
$Q_0:=Q\oplus Q$ on $H_0$.  Choose $R_0>0$ so large that
\[
  \overline{\mathcal C^0}\subset(-R_0,R_0)\times[0,d],
  \qquad
  \rho=\rho^-\text{ for }x<-R_0,
  \qquad
  \rho=\rho^+\text{ for }x>R_0.
\]
The last two equalities hold by construction of the two periodic half-guides.
All norms below may be compared with their unweighted counterparts because of
the common positive upper and lower bounds on the coefficients.
For a scalar operator $T$, write $\rho_T=\rho$ when $T=A$ and
$\rho_T=\rho^\star$ when $T=A^\star$; formulas for $A_0$ are understood
componentwise.

\begin{lemma}[Local compactness and cutoff commutators]
\label{lem:appendix-local-compactness}
Let $T$ be any one of $A,A^-,A^+$, or $A_0$, acting in its corresponding
Hilbert space $H_T$, and let $Q_T$ be the corresponding position operator.
Then:
\begin{enumerate}[label=\textup{(\roman*)},leftmargin=*]
  \item for every bounded interval $K\subset\R$,
  \[
    \mathbf 1_K(Q_T)(T+\ii)^{-1}:H_T\longrightarrow H_T
    \quad\text{is compact};
  \]
  \item if $\zeta\in C_c^\infty(\R;[0,1])$ satisfies
  $\zeta=1$ on $[-1,1]$ and $\zeta=0$ outside $[-2,2]$, and
  $\zeta_R(x):=\zeta(x/R)$, then
  \begin{equation}
    \bigl\|[T,\zeta_R(Q_T)](T+\ii)^{-1}\bigr\|_{\mathcal B(H_T)}
    \longrightarrow0
    \qquad(R\to\infty).
    \label{eq:appendix-commutator-limit}
  \end{equation}
\end{enumerate}
\end{lemma}
\begin{proof}
We first treat a scalar component.  Its closed form is
\[
  a_T(u,v)=\int_B\nabla u\cdot\nabla\overline v\,\dd x\dd y,
  \qquad D(a_T)=H^1_\beta(B).
\]
If $u=(T+\ii)^{-1}f$, nonnegativity and the variational identity give
\[
  \|u\|_{H_T}\leq\|f\|_{H_T},
  \qquad
  \|\nabla u\|_{L^2(B)}^2
  =\operatorname{Re}\langle f,u\rangle_{H_T}
  \leq\|f\|_{H_T}^2.
\]
Thus the resolvent maps boundedly into $H^1_\beta(B)$.  Restriction to
$K\times(0,d)$ followed by the Rellich embedding
$H^1(K\times(0,d))\hookrightarrow L^2(K\times(0,d))$ is compact.  The
coefficient bounds make the weighted and unweighted local $L^2$ norms
equivalent, proving (i).  The direct-sum case follows componentwise.

Multiplication by a smooth function of $x$ preserves the operator domain:
for $u\in H^1(\Delta,B)$,
\[
  \Delta(\eta u)=\eta\Delta u+2\eta'\partial_xu+\eta''u\in L^2(B),
\]
and the horizontal $\beta$-quasiperiodic trace conditions are unchanged.
Consequently, on $D(T)$,
\begin{equation}
  [T,\zeta_R]u
  =-\frac1{\rho_T}
    \left(2\zeta_R'\partial_xu+\zeta_R''u\right),
  \label{eq:appendix-commutator-formula}
\end{equation}
with the formula interpreted componentwise for $A_0$.  Since
$\|\zeta_R'\|_\infty=O(R^{-1})$ and
$\|\zeta_R''\|_\infty=O(R^{-2})$, the preceding resolvent estimate yields
\[
  \bigl\|[T,\zeta_R](T+\ii)^{-1}f\bigr\|_{H_T}
  \leq C\left(R^{-1}+R^{-2}\right)\|f\|_{H_T}.
\]
This proves (ii).
\end{proof}

\begin{lemma}[Zhislin criterion on the strip]
\label{lem:appendix-Zhislin}
For any operator $T$ in
\Cref{lem:appendix-local-compactness} and any $\lambda\in\R$, the following
are equivalent:
\begin{enumerate}[label=\textup{(\roman*)},leftmargin=*]
  \item $\lambda\in\sigma_{\mathrm{ess}}(T)$;
  \item there exists $u_m\in D(T)$ such that
  \begin{equation}
    \|u_m\|_{H_T}=1,\qquad
    \mathbf 1_{[-m,m]}(Q_T)u_m=0,\qquad
    \|(T-\lambda)u_m\|_{H_T}\longrightarrow0.
    \label{eq:appendix-Zhislin-sequence}
  \end{equation}
\end{enumerate}
For $T=A_0$, the support condition is imposed on both components.
\end{lemma}
\begin{proof}
A sequence satisfying (ii) converges weakly to zero: this is immediate first
against compactly supported $L^2$ functions and then against arbitrary
functions by density.  It is therefore a singular Weyl sequence, which proves
(i).

Conversely, let $v_j$ be a singular Weyl sequence for $T$ and $\lambda$.
For every fixed $R$, local compactness gives
\[
  \zeta_R(Q_T)v_j
  =\zeta_R(Q_T)(T+\ii)^{-1}(T+\ii)v_j
  \longrightarrow0
  \quad\text{strongly as }j\to\infty.
\]
Indeed,
$(T+\ii)v_j=(T-\lambda)v_j+(\lambda+\ii)v_j$ converges weakly to zero and
is bounded.  Choose a subsequence, still denoted $v_m$, such that
\[
  \|\zeta_m(Q_T)v_m\|_{H_T}\leq m^{-1},\qquad
  \|(T-\lambda)v_m\|_{H_T}\leq m^{-1}.
\]
Set $w_m=(1-\zeta_m(Q_T))v_m$.  Then $w_m$ vanishes for $|x|\leq m$ and
$\|w_m\|_{H_T}\to1$.  Moreover,
\[
  (T-\lambda)w_m
  =(1-\zeta_m)(T-\lambda)v_m-[T,\zeta_m]v_m.
\]
The first term tends to zero.  For the second, write
\[
  [T,\zeta_m]v_m
  =[T,\zeta_m](T+\ii)^{-1}(T+\ii)v_m.
\]
The first factor tends to zero by
\eqref{eq:appendix-commutator-limit}, while the second is bounded.  Normalizing
$w_m$ gives \eqref{eq:appendix-Zhislin-sequence}.  This is the strip analogue
of the Zhislin construction in
\cite[Lemmas~3.11--3.12]{DeNittisEtAl2021}.
\end{proof}

\begin{lemma}[Periodic Weyl sequences placed in a prescribed end]
\label{lem:appendix-one-sided-Weyl}
Let $\star\in\{-,+\}$ and
$\lambda\in\sigma(A^\star)$.  There are normalized
$v_m^\star\in D(A^\star)$ such that
\[
  \|(A^\star-\lambda)v_m^\star\|_{H^\star}\longrightarrow0
\]
and
\[
  \operatorname{supp}v_m^+\subset[m,\infty)\times[0,d],
  \qquad
  \operatorname{supp}v_m^-\subset(-\infty,-m]\times[0,d].
\]
\end{lemma}
\begin{proof}
We give the right-end construction.  By the Weyl criterion for the spectrum,
choose $f_m\in D(A^+)$ with
\[
  \|f_m\|_{H^+}=1,\qquad
  \|(A^+-\lambda)f_m\|_{H^+}\leq m^{-1}.
\]
Since $A^+$ is nonnegative and associated with the gradient form,
\[
  \|\nabla f_m\|_{L^2(B)}^2
  =\langle A^+f_m,f_m\rangle_{H^+}
  \leq\lambda+m^{-1}
\]
for all sufficiently large $m$.  The periodic translation
\[
  (\tau_q^+f)(x,y):=f(x-qL^+,y),\qquad q\in\Z,
\]
is unitary in $H^+$ and commutes with $A^+$.  Choose $q_m$ so large that
\[
  \bigl\|\mathbf 1_{(-\infty,2m]}(Q)\tau_{q_m}^+f_m
  \bigr\|_{H^+}\leq m^{-1}.
\]
Let $\eta\in C^\infty(\R;[0,1])$ satisfy
$\eta(t)=0$ for $t\leq0$ and $\eta(t)=1$ for $t\geq1$, and set
$\eta_m^+(x):=\eta((x-m)/m)$.  Then
\[
  g_m:=\eta_m^+\tau_{q_m}^+f_m
\]
is supported in $[m,\infty)\times[0,d]$ and
$\|g_m\|_{H^+}\to1$.  The domain-preservation calculation in
\Cref{lem:appendix-local-compactness} gives $g_m\in D(A^+)$, and
\[
  (A^+-\lambda)g_m
  =\eta_m^+\tau_{q_m}^+(A^+-\lambda)f_m
   +[A^+,\eta_m^+]\tau_{q_m}^+f_m.
\]
The first term tends to zero.  Formula
\eqref{eq:appendix-commutator-formula}, the uniform $H^1$ bound, and
$\|\partial_x^r\eta_m^+\|_\infty=O(m^{-r})$ for $r=1,2$ show that the second
term also tends to zero.  Normalize $g_m$ to obtain $v_m^+$.  Translating by
negative integer multiples of $L^-$ and using
$\eta_m^-(x):=\eta((-x-m)/m)$ gives the left-end sequence.  This is the
second-order strip counterpart of the translation-and-cutoff construction in
the proof of
\cite[Proposition~3.10]{DeNittisEtAl2021}.
\end{proof}

\begin{proof}[Proof of \Cref{thm:two-ended-essential-spectrum}]
Floquet--Bloch decomposition applied separately to the two periodic
extensions gives, by \eqref{eq:background-bands},
\begin{equation}
  \sigma_{\mathrm{ess}}(A_0)
  =\sigma_{\mathrm{ess}}(A^-)
   \cup\sigma_{\mathrm{ess}}(A^+)
  =\sigma(A^-)\cup\sigma(A^+)
  =\Sigma_-(\beta)\cup\Sigma_+(\beta).
  \label{eq:appendix-free-essential-spectrum}
\end{equation}
Here the first equality is the elementary essential-spectrum formula for a
finite orthogonal direct sum, while the second is the periodic-spectrum result
of \cite[Proposition~3.1, especially equations~(3.3)--(3.4)]{Fliss2013}.

We first prove
$\sigma_{\mathrm{ess}}(A)\subset\sigma_{\mathrm{ess}}(A_0)$.
Choose $j_-,j_+\in C^\infty(\R;[0,1])$ such that
\[
\begin{array}{lll}
  j_-(x)=1&\text{for }x\leq-2R_0,&
  j_-(x)=0\text{ for }x\geq-R_0,\\
  j_+(x)=0&\text{for }x\leq R_0,&
  j_+(x)=1\text{ for }x\geq2R_0.
\end{array}
\]
Let $\lambda\in\sigma_{\mathrm{ess}}(A)$ and take the Zhislin sequence
$u_m$ furnished by \Cref{lem:appendix-Zhislin}.  For $m>2R_0$,
\[
  u_m=j_-u_m+j_+u_m,
\]
and the two terms have disjoint supports.  Define
\[
  u_m^0:=c_m(j_-u_m,j_+u_m)\in D(A_0),
  \qquad
  c_m:=\|(j_-u_m,j_+u_m)\|_{H_0}^{-1}.
\]
On the support of $j_-u_m$ the weights $\rho$ and $\rho^-$ agree, and on the
support of $j_+u_m$ the weights $\rho$ and $\rho^+$ agree.  Hence $c_m=1$ for
all sufficiently large $m$.  The derivatives of $j_\pm$ are supported in the
fixed set $[-2R_0,2R_0]$, where $u_m$ vanishes.  Therefore the commutator
terms vanish and
\[
\begin{aligned}
  \|(A_0-\lambda)u_m^0\|_{H_0}^2
  &=\|(A^--\lambda)j_-u_m\|_{H^-}^2
    +\|(A^+-\lambda)j_+u_m\|_{H^+}^2\\
  &\leq C\|(A-\lambda)u_m\|_H^2\longrightarrow0.
\end{aligned}
\]
The sequence $u_m^0$ still vanishes in $|x|\leq m$.  By
\Cref{lem:appendix-Zhislin},
$\lambda\in\sigma_{\mathrm{ess}}(A_0)$.

For the reverse inclusion, let
$\lambda\in\sigma(A^+)$.  By
\Cref{lem:appendix-one-sided-Weyl}, there is a normalized approximate
eigenvector sequence $v_m^+$ supported in
$[m,\infty)\times[0,d]$.  Once $m>R_0$, the two weights and the two
differential realizations agree on this support.  Thus
\[
  \|v_m^+\|_H=1,\qquad
  \|(A-\lambda)v_m^+\|_H
  =\|(A^+-\lambda)v_m^+\|_{H^+}\longrightarrow0.
\]
Because the supports escape to the right, $v_m^+\rightharpoonup0$ in $H$.
It is a singular Weyl sequence for $A$, and hence
$\lambda\in\sigma_{\mathrm{ess}}(A)$.  The left-end construction gives
$\sigma(A^-)\subset\sigma_{\mathrm{ess}}(A)$ in the same way.  Consequently,
\[
  \sigma(A^-)\cup\sigma(A^+)
  \subset\sigma_{\mathrm{ess}}(A)
  \subset\sigma_{\mathrm{ess}}(A_0).
\]
Combining this with \eqref{eq:appendix-free-essential-spectrum} proves
\eqref{eq:projected-spectrum}.
\end{proof}


% Full title: Corrected Representation Identity
\section{Corrected representation identity}
\label{app:corrected-identity}

The wavenumber assignment in the complementary-field argument is recorded
here to prevent a common mismatch.  If $w^{\mathrm e}$ solves the background
equation $(\Delta+k^2)w^{\mathrm e}=0$ and $v^{\mathrm i}$ is a solution of
the complementary interior-$k$ equation, the quasiperiodic extinction pair is
\begin{equation}
  -S^{\mathrm{qp}}_{k,\beta}
  [\gamma_N^{\mathrm i}v^{\mathrm i}]
  +D^{\mathrm{qp}}_{k,\beta}
  [\gamma_D^{\mathrm i}v^{\mathrm i}].
  \label{eq:corrected-exterior-pair}
\end{equation}
Both operators in \eqref{eq:corrected-exterior-pair} use the background
wavenumber $k$.  The inclusion representation uses the free-space pair
\begin{equation}
  -S_{k^0}[\gamma_N^{\mathrm i}w^{\mathrm i}]
  +D_{k^0}[\gamma_D^{\mathrm i}w^{\mathrm i}],
  \qquad k^0=n^0k.
  \label{eq:corrected-interior-pair}
\end{equation}
Mixing $D^{\mathrm{qp}}_{k^0,\beta}$ into
\eqref{eq:corrected-exterior-pair} invalidates the exterior $k$-Helmholtz
extinction identity.  Equations \eqref{eq:corrected-exterior-pair} and
\eqref{eq:corrected-interior-pair} agree with the wavenumber separation used
in \eqref{eq:center-reconstruction} and with the architecture of
\cite[Appendix~A]{BarnettGreengard2010}.


\printbibliography

\end{document}
